\documentclass[11pt]{amsart}

\usepackage[T1]{fontenc}
\usepackage{amsmath,amssymb,amsthm}
\usepackage{mathtools}
\usepackage{booktabs}
\usepackage{array}
\usepackage{longtable}
\usepackage{enumitem}
\usepackage[margin=1.15in]{geometry}
\usepackage[colorlinks=true,linkcolor=blue!50!black,citecolor=green!40!black,urlcolor=blue!50!black]{hyperref}
\usepackage{xcolor}

%-------------------------------------------------------------------------------
% Theorem environments
%-------------------------------------------------------------------------------
\newtheorem{mainthm}{Theorem}
\renewcommand{\themainthm}{\Alph{mainthm}}
\newtheorem{thmB}{Theorem}
\renewcommand{\thethmB}{B\arabic{thmB}}
\newtheorem{corB}[thmB]{Corollary}
\newtheorem*{thmBsummary}{Theorem B}

\newtheorem{theorem}{Theorem}[section]
\newtheorem{proposition}[theorem]{Proposition}
\newtheorem{lemma}[theorem]{Lemma}
\newtheorem{corollary}[theorem]{Corollary}
\newtheorem{conjecture}[theorem]{Conjecture}
\theoremstyle{definition}
\newtheorem{definition}[theorem]{Definition}
\theoremstyle{remark}
\newtheorem{remark}[theorem]{Remark}

%-------------------------------------------------------------------------------
% Macros
%-------------------------------------------------------------------------------
\newcommand{\Q}{\mathbb{Q}}
\newcommand{\Z}{\mathbb{Z}}
\newcommand{\C}{\mathbb{C}}
\newcommand{\F}{\mathbb{F}}
\newcommand{\PP}{\mathbb{P}}
\newcommand{\cO}{\mathcal{O}}
\newcommand{\fp}{\mathfrak{p}}
\newcommand{\fP}{\mathfrak{P}}
\DeclareMathOperator{\Gal}{Gal}
\DeclareMathOperator{\GL}{GL}
\DeclareMathOperator{\SL}{SL}
\DeclareMathOperator{\PGL}{PGL}
\DeclareMathOperator{\PSL}{PSL}
\DeclareMathOperator{\Li}{Li}
\DeclareMathOperator{\tr}{tr}
\DeclareMathOperator{\Frob}{Frob}
\DeclareMathOperator{\Ind}{Ind}
\DeclareMathOperator{\Cl}{Cl}
\DeclareMathOperator{\Nm}{N}
\newcommand{\Inv}{\mathrm{Inv}}
\newcommand{\Zt}{\mathrm{Z}}       % hypothesis Z(theta)
\newcommand{\Nor}{\mathcal{N}}     % normaliser of a torus
\newcommand{\Nout}{\Nor^{-}}
\newcommand{\inert}{\mathrm{inert}}
\newcommand{\cA}{\mathcal{A}}      % certificate main-term coefficient
\newcommand{\cE}{\mathcal{E}}      % certificate error coefficient
\newcommand{\cG}{\mathcal{G}}      % Gal(Omega/Q)

\emergencystretch=1.5em

\begin{document}

\title[Vanishing of Fourier coefficients and involutions in $\PGL_2(\F_\ell)$]{Vanishing of Fourier coefficients from
a uniform zero-free half-plane, and an uncertainty principle for involutions in $\PGL_2(\F_\ell)$}

\author{Morgan Blem}
\urladdr{https://github.com/SlackEight/tau-vanishing}

\date{14 October 2026. Preprint, version 1.1.2; not refereed. Prepared with substantial AI assistance; see
Section~\ref{sec:provenance}. Lean code and computations: \url{https://github.com/SlackEight/tau-vanishing}}

\begin{abstract}
Write $\Zt(\theta)$ for the statement that every Dedekind zeta function is nonzero in $\operatorname{Re}s>\theta$,
uniformly in the field. Cureton and Thorner state $\Zt(1-1/6000)$ in a preprint; their Lean development (proofs
generated by AxiomProver) proves $\Zt(31/32)$, a constant that is not in the preprint (v1). We record the transfer of the Serre--Murty--Murty--Saradha--Zywina argument to this setting:
for a non-CM newform $f$ of even weight $k\ge2$ with integer coefficients, $\Zt(\theta)$ implies
$\#\{p\le x:a_f(p)=0\}\ll_f x^{(1+\theta)/2}$; given the Lean theorem, this is $x^{63/64}$ for Ramanujan's $\tau$.
This implication is verified in Lean~4 in the form $x^{(1+\theta)/2+\varepsilon}$, conditional on an explicit
hypothesis encoding the mod-$\ell$ Galois representations and their large image. The main new ingredient is an exact
inequality for the involutions in $\PGL_2(\F_\ell)$: every nonnegative class-function majorant $\varphi$ of the involutions
satisfies $\mu(\varphi)\lambda_A(\varphi)\ge\kappa_\ell\ge1-1/\ell$, where $\mu$ is the mean and $\lambda_A$ is
Bella\"iche's Littlewood norm. With a companion inequality, it gives lower bounds for three explicitly defined error
functionals, attached to single-$\ell$ majorants, to single-extension Chebotarev certificates (over fields whose
Galois group maps onto the product of the groups $\PGL_2(\F_\ell)$ with central kernel), and to per-prime Selberg sieves with pairwise error accounting (a class that contains Bella\"iche's large sieve). None of these
functionals can certify an exponent below $(1+\theta)/2$.
\end{abstract}

\maketitle

\tableofcontents

%===============================================================================
\section{Introduction}\label{sec:intro}
%===============================================================================

\subsection{The hypothesis and the inputs}

For a number field $K$ and $\tfrac12\le\theta<1$ we write
\[
  \Zt_K(\theta):\qquad \zeta_K(s)\neq 0\quad\text{for } \operatorname{Re}s>\theta,
\]
where $\zeta_K$ denotes the meromorphic continuation of the Dedekind zeta function. We write $\Zt(\theta)$ for the
statement that $\Zt_K(\theta)$ holds for \emph{every} number field $K$, with the same $\theta$. All results of this
note are implications from $\Zt(\theta)$, or from $\Zt_K(\theta)$ for the fields named in each statement. For the main
application the uniformity in $K$ is essential, because the fields involved vary with the counting parameter $x$
(Remark~\ref{rem:uniformity}).

Two recent sources supply $\Zt(\theta)$ for a specific $\theta$. We state the corollaries that depend on them
separately, as conditional on these inputs.
\begin{enumerate}[label=(I\arabic*),leftmargin=*]
\item \emph{The Cureton--Thorner preprint} \cite{CT}, Theorem~1.1: \emph{if $K$ is any number field, then
  $\zeta_K(s)\neq0$ when $\operatorname{Re}s>1-1/6000$.} This is $\Zt(1-1/6000)$. The constant is absolute, and there is
  no height cutoff, no exceptional-zero exemption and no field-dependent threshold. Its Corollary~1.2 gives the same
  half-plane for every finite-order Hecke $L$-function. The preprint is unrefereed; according to its authors, its
  proofs were generated and verified in Lean by AxiomProver. It is posted on Hexagon (\url{https://hexagonmath.org}),
  a repository for research works in the mathematical sciences operated by the Hexagon Mathematics Foundation.
  Submissions to Hexagon are screened (for compilation errors, scholarly merit and offensive content) but not
  refereed, and each posted version keeps a permanent identifier, here \texttt{hexagon:2610.00200v1}.
\item \emph{The Lean development of Cureton--Thorner} \cite{CT-Lean} (D.~Cureton, J.~Thorner; proofs generated by
  AxiomProver), repository \texttt{AxiomMath/DedekindZetaZeroFree}, commit \texttt{09afbb4c} (full hash in
  \cite{CT-Lean}). It proves
\begin{verbatim}
theorem main (K : Type) [Field K] [NumberField K] :
  exists f : C -> C, DifferentiableOn C f {1}^c /\
    (forall s, 1 < s.re -> f s = NumberField.dedekindZeta K s) /\
    forall s, 31/32 < s.re -> s != 1 -> f s != 0
\end{verbatim}
  (displayed here in ASCII), together with the analogous statement for finite-order Hecke $L$-functions. Here
  \texttt{NumberField.dedekindZeta} is Mathlib's Dedekind zeta function, which is not shadowed in the repository.
  Since $\C\setminus\{1\}$ is connected, $f$ is the meromorphic continuation of $\zeta_K$, so the theorem asserts
  $\Zt(31/32)$. The constant $31/32$ appears so far only in this Lean development, not in the preprint \cite{CT}
  (v1), and as of this writing no human-readable proof of it is public. Its correctness rests on the Lean kernel,
  Mathlib's definitions and the toolchain. The $31/32$ commit (\texttt{201bc40}) is dated 13 October 2026, 09:09 UTC,
  and the Hecke corollary commit (\texttt{09afbb4c}) 10:10 UTC the same day. We know of no archival copy of this
  commit (Software Heritage or Zenodo) as of 14 October 2026; its intrinsic identifier is
  \texttt{swh:1:rev:09afbb4ce6570c1615d58509dc197fc9f0956b7d}, which resolves once the repository is archived.

  We built the full closure (Lean \texttt{v4.35.0-rc3}); \texttt{\#print axioms} for \texttt{main} reports only
  \texttt{propext}, \texttt{Classical.choice} and \texttt{Quot.sound}. A third-party repository \cite{Murphy}
  reports that the proof replays, but that the Comparator challenge shipped with \cite{CT-Lean} is not accepted as
  shipped: an instance path (\texttt{instCommCStarAlgebraComplex}) makes the challenge and solution statements
  syntactically different though definitionally equal, and a one-line change to the challenge makes it pass. We ran
  Comparator (\texttt{lake comparator}, Lean's default kernel) on the challenge and \texttt{comparator.json} shipped
  with \cite{CT-Lean} at \texttt{09afbb4c}, unmodified, on 14 October 2026: it reported ``Lean default kernel accepts
  the solution'' and ``Your solution is okay!'' (exit code~0) for \texttt{main} and \texttt{heckeCorollary}. So the
  reported mismatch did not reproduce in our run. We have not determined which version of the challenge that report
  concerns; the history of the challenge file includes a commit ``fix comparator challenge imports''
  (\texttt{896b677}, 12 October 2026). Our own Comparator run (Appendix~\ref{sec:lean}) checks our
  theorems against our own challenge file, and the closure it replays through the kernel includes the proof of
  \texttt{DedekindZetaZeroFree.main}.
\end{enumerate}

\emph{Context.} Cureton--Thorner follow the strategy of OpenAI's quasi-Riemann hypothesis \cite{OAI-QRH} (every
Dirichlet $L$-function, and $\zeta$ itself, is nonzero for $\operatorname{Re}s>7/8$; the statement for $\zeta$ is
formalised in \cite{OAI-Lean}), through OpenAI's primitive-roots paper \cite{OAI-PR}. That work does not give
$\Zt_K(\theta)$ for general $K$, and no result here uses it. Hu \cite{Hu} (\texttt{hexagon:2610.00216v1}) claims
nonvanishing for $\operatorname{Re}s>7/8$ for every Artin $L$-function over every number field; that paper is
unrefereed, AI-generated and not formalised, and we mention it only to illustrate the value $\theta=7/8$.

%-------------------------------------------------------------------------------
\subsection{Results}
%-------------------------------------------------------------------------------

\begin{itemize}[leftmargin=*]
\item \textbf{Theorem A} (\S\ref{sec:A}). Let $f$ be a non-CM newform of even weight $k\ge 2$, trivial character and
  integer coefficients. Then $\Zt(\theta)$ implies
  \[\#\{p\le x:\ a_f(p)=0\}\ll_{f,\theta} x^{(1+\theta)/2}.\]
  Given (I2), $\#\{p\le x:\tau(p)=0\}\ll x^{63/64}$; given (I1), $\ll x^{1-1/12000}$ (Corollaries~\ref{cor:A-I2}
  and~\ref{cor:A-I1}). In weight~$2$ Elkies' unconditional $x^{3/4}$ is better, so the statement is of interest for
  $k\ge4$. The proof is the standard GRH argument with $\theta$ in place of $\tfrac12$, run in a single quadratic
  extension $M_\ell/E_\ell$ with $\ell\asymp x^{(1-\theta)/2}/\log x$. The implication is verified in Lean~4 in the
  form $x^{(1+\theta)/2+\varepsilon}$, conditional on an explicit hypothesis encoding Deligne's mod-$\ell$
  representations and their large image (Appendix~\ref{sec:lean}).
\end{itemize}

Theorem~B concerns three classes of upper-bound arguments, each with an explicitly postulated error functional
(Definition~\ref{def:classes}): \emph{single-$\ell$ majorants} $\varphi$ on $\PGL_2(\F_\ell)$, with certified bound
$B_{\mathrm{maj}}=\mu(\varphi)\Li(x)+c\,\lambda(\varphi)x^\theta w(x)$; \emph{single-extension certificates} $(L/K,C,m)$
built from mod-$\ell$ representations, with $B_{\mathrm{cert}}=\cA\Li(x)+c\,\cE x^\theta\log x$, where $\cA$ and
$\cE$ are the main-term and Lagarias--Odlyzko error coefficients; and \emph{per-prime Selberg sieves} with weights
$y$, with $B_{\mathrm{sieve}}=\Li(x)\sum y_m^2+c\,x^\theta w(x)\sum_{m,n}|y_my_n|E(m,n)$. In each class the test
object is required to be valid on the whole set of elements with trace zero, not only at the Frobenius classes of the
primes with $a(p)=0$.

\begin{thmBsummary}[\S\ref{sec:B}: Theorem~\ref{thm:B1}, Corollary~\ref{cor:B2}, Theorem~\ref{thm:B3} and
Proposition~\ref{prop:sieve}]
Let $\ell\ge5$ be prime and let $\Inv_\ell$ be the set of involutions of $\PGL_2(\F_\ell)$.
\begin{enumerate}[label=(\roman*)]
\item Every class-function majorant $\varphi$ of $\Inv_\ell$ satisfies $\mu(\varphi)\lambda_A(\varphi)\ge\kappa_\ell
  =\ell(\ell^2+1)/((\ell+1)(\ell^2-1))\ge1-1/\ell$, with equality for $\varphi=1_{\Inv}$.
\item Let $f$ be as in Theorem~A and $\tfrac12\le\theta<1$. For every choice of data in the three classes of
  Definition~\ref{def:classes} (with the additional hypotheses of Theorem~\ref{thm:B3} and
Proposition~\ref{prop:sieve}), the certified bound is
  $\ge x^{(1+\theta)/2-o(1)}$ as $x\to\infty$; explicit lower bounds are given in Corollary~\ref{cor:B2},
  Theorem~\ref{thm:B3} and Proposition~\ref{prop:sieve}.
\end{enumerate}
\end{thmBsummary}

Part (ii) is a statement about the three functionals, and about nothing else. Methods whose error terms have a
different shape are not covered; \S\ref{subsec:B-scope} lists natural variants that are not, among them smoothing,
Hecke-shape errors, information modulo $\ell^m$, sign-changing or cancelling sieve weights, and nonlinear selections
of $\ell$. The proof of Theorem~A is of the first kind, and Bella\"iche's large sieve \cite[Thm.~14]{Bel16} is of the
third. The arXiv preprint \cite{Bel13} claimed $x^{2/3+\varepsilon}$ for $a=0$ (under GRH and Artin's conjecture);
the published version \cite{Bel16} no longer contains that bound and acknowledges two errors in the sieve argument of
an earlier version, one found by Zywina and a second found later by L.~Devin and an anonymous referee
\cite[introduction and Rem.~18(i)]{Bel16}. Its corrected sieve bound gives $x^{3/4+\varepsilon}$ for
$a=0$, and exactly $x^{(1+\theta)/2+\varepsilon}$ under $\Zt(\theta)$ (\S\ref{subsec:B-scope}). At $\theta=\frac12$
the lower bounds of Theorem~B are $x^{3/4-o(1)}$, the exponent of Serre \cite{Serre81}, of Zywina \cite{Zyw15} (whose
proof follows Murty--Murty--Saradha \cite{MMS}) and of \cite[Cor.~15]{Bel16}. This is consistent with the fact that the exponent $3/4$ for
$a=0$ under GRH has not been improved without further hypotheses, but Theorem~B does not explain that fact: it covers
only the three classes defined here.

\S\ref{sec:further} records three further consequences of $\Zt(\theta)$, each an immediate consequence or an
instantiation of published work.

\begin{table}[h]
\centering
\small
\begin{tabular}{@{}p{0.2\textwidth}p{0.3\textwidth}p{0.42\textwidth}@{}}
\toprule
item & content & status (author's assessment) \\
\midrule
Thm~A & $\Zt(\theta)\Rightarrow\#\{a_f(p)=0\}\ll x^{(1+\theta)/2}$ & immediate consequence of $\Zt(\theta)$: transfer
  of the Serre/MMS/Zywina argument (\S\ref{subsec:novelty}); formalised in Lean with $+\varepsilon$, given H1
  (Appendix~\ref{sec:lean}) \\
Thm~\ref{thm:B1} & exact uncertainty inequality on $\PGL_2(\F_\ell)$ & not found in the literature; short \\
Cor.~\ref{cor:B2}, Thm.~\ref{thm:B3}, Prop.~\ref{prop:sieve} & lower bounds for three error functionals &
  not found in the literature; modest (uncertainty inequalities plus AM--GM); Thm.~\ref{thm:B3} uses a finite
  computation for $\ell<61$ \\
\midrule
Prop.~\ref{prop:C} & EV $\ell$-torsion bound, every $K$ & immediate consequence; also \cite[Cor.~1.4]{FSW} \\
Prop.~\ref{prop:D} & Koblitz upper bound, elliptic Brun & instantiation (cf.\ \cite[footnote~1]{Zyw08}) \\
Prop.~\ref{prop:E} & CM degree $\le216$ ($\le252$ from (I1)) & short deduction from Lee--Kwon; uses their tables \\
\bottomrule
\end{tabular}
\caption{Status at a glance. The assessments are the author's. The note has not been refereed, and the author's own
verification of the proofs is in progress (\S\ref{sec:provenance}).}
\label{tab:status}
\end{table}

\subsection{Relation to prior work, and what is new}\label{subsec:novelty}

\emph{Theorem A is not new in substance.} It is what one obtains by inserting $\Zt(\theta)$ into the standard GRH
proofs of the bound $x^{3/4}$ (Serre \cite[Cor.~2 to Thm.~15]{Serre81}; Zywina \cite[\S4, Rem.~1.5]{Zyw15},
whose proof follows Murty--Murty--Saradha \cite{MMS}), and we regard it as an immediate consequence of the uniform half-plane, of the same status as
Proposition~\ref{prop:C}. In particular:
\begin{itemize}[leftmargin=*]
\item Murty, Murty and Pujahari \cite[\S7(b), p.~12]{MMP} write, about $\#\{p\le n:\tau(p)=0\}$: ``One can estimate
  the number of such primes $\le n$ to be $O(n^\theta)$ with $\theta<1$, under a quasi-generalized Riemann hypothesis.
  This is easy to derive following a mild variation in the argument of section~4 of [MMS].'' Zywina's proof of $x^{3/4}(\log x)^{-1/2}$ under GRH, described
  next, follows the argument of \cite{MMS}. (MMP's own eq.~(7), $\pi(n)/\ell+O(\ell^3n^{1/2}\log n)$ in
  $\GL_2(\F_\ell)$, gives $n^{7/8}$ under GRH, and $n^{(3+\theta)/4}$ with $n^\theta$ in place of $n^{1/2}$.)
\item Zywina's proof for $a=0$ \cite[\S4, last paragraph]{Zyw15} works in a single cyclic extension $L^H/L^B$ of
  degree $\ell-1$ over a base of degree $\ell+1$, with $|C''|=1$. Under $\Zt(\theta)$ and with a sharp cutoff, the
  same lines give $x/(\ell\log x)+\ell x^\theta\log x$, hence the exponent $(1+\theta)/2$. His remaining use of GRH, the
  selection of $\ell$ \cite[Lemma~4.4]{MMS}, has an unconditional substitute, V.~K.~Murty's inequality quoted as
  \cite[(9.6)]{TZ}. That substitute is not lossless: (9.6) adds a secondary term
  $(\pi_f(x,a)\,x\log y/y)^{1/2}$, which after solving for $\pi_f(x,a)$ costs $\asymp x\log y/y$, that is
  $x^{(1+\theta)/2}(\log x)^2$ for $y\asymp x^{(1-\theta)/2}/\log x$. The exponent is unchanged; the power of $\log x$
  is worse. (Thorner and Zaman question how V.~K.~Murty's final bound follows from (9.6) and (9.7) \cite[\S9.1]{TZ},
  not (9.6) itself.)
\item Serre \cite[\S7.3(iii)]{Serre81} already passes to the projective image, where the trace-zero locus becomes the
  set of (regular) elements of order~$2$. The quadratic extension $M_\ell/E_\ell=L'^{T}/L'^{\Nor}$ of Theorem~A is the
  mod-$\ell$ version of this observation.
\item Cojocaru--Meyer \cite[Lemma~11, (43)]{CM} state $x^{(2+\theta)/3}(\log x)^{-1/3}$ under a $\theta$-hypothesis
  for elliptic curves, and say (Rem.~12) that much better conditional bounds are known.
\end{itemize}
What this note adds to Theorem~A is: the check that the uniformity of (I1) and (I2) is what the argument needs
(Remark~\ref{rem:uniformity}); a formulation in $\PGL_2(\F_\ell)$ with $\Nor/T$ that captures split and nonsplit
involutions at once, so that no selection lemma over $\ell$ is needed (this is the finite-level form of Serre's
\cite[\S7.3(iii)]{Serre81}); and a Lean formalisation of the implication (Appendix~\ref{sec:lean}).

\emph{Theorem B.} We did not find Theorem~\ref{thm:B1}, Lemma~\ref{lem:S1} or the lower bounds of
Corollary~\ref{cor:B2}, Theorem~\ref{thm:B3} and Proposition~\ref{prop:sieve} in the literature. They are modest: a
finite-group linear-programming fact, and its consequences for explicitly postulated error functionals, which make
precise an expectation that is informally present in the literature. The informal message that Chebotarev methods stop
at $x^{3/4}$ under GRH is in Zywina \cite[\S1.3]{Zyw15}; neither he nor anyone we know of proves a statement of this
kind. Proposition~\ref{prop:sieve} makes precise, for the functional $B_{\mathrm{sieve}}$, the remark of Zywina
\cite[Rem.~1.7(iii)]{Zyw08} that it is better to sieve by a single well-chosen prime than by many. The closest
structural analogue of Theorem~\ref{thm:B1} that we know is Bella\"iche's Remark~16 in \cite{Bel16}
(\S\ref{subsec:B-scope}).

\subsection{Organisation}\label{subsec:organisation}
\S\ref{sec:cheb} states the Chebotarev input, \S\ref{sec:A} proves Theorem~A and \S\ref{sec:B} proves Theorem~B.
Theorem~\ref{thm:B1} is a statement about finite groups and depends on nothing else; the rest of \S\ref{sec:B} uses
from \S\ref{sec:A} only the large-image set-up and the decomposition dictionary. \S\ref{sec:further} records further
consequences of $\Zt(\theta)$, which are independent of Theorems~A and~B and of each other. \S\ref{sec:computations}
lists the computations, \S\ref{sec:provenance} describes how the note was prepared, Appendix~\ref{sec:lean}
describes the Lean formalisation, and Appendix~\ref{sec:deps} tabulates what each statement depends on and how it is
checked.

\subsection*{Notation}
$p$ always denotes a rational prime, and $N$ the level of a newform. $\Li(x)=\int_2^x dt/\log t$. For a number field $L$:
$n_L=[L:\Q]$, $d_L=|\mathrm{disc}(L)|$, and $D_K=d_K$. $\omega(m)$ is the number of distinct prime factors of $m$.
Implied constants are effective unless stated otherwise.

%===============================================================================
\section{Chebotarev under a zero-free half-plane}\label{sec:cheb}
%===============================================================================

\begin{proposition}[Artin $L$-functions from $\zeta_L$]\label{prop:artin}
Let $L/K$ be a Galois extension of number fields and assume $\Zt_L(\theta)$. Then every Artin $L$-function
$L(s,\rho,L/K)$ is holomorphic and nonvanishing in $\operatorname{Re}s>\theta$, except for a pole at $s=1$ of order
$\langle\rho,1\rangle$.
\end{proposition}

\begin{proof}
By Brauer's induction theorem, $L(s,\rho)=\prod_iL(s,\psi_i,L/E_i)^{n_i}$ with $n_i\in\Z$ and $\psi_i$ linear
characters of $H_i=\Gal(L/E_i)$. Fix $i$ and put $M_i=L^{\ker\psi_i}$. Then $\zeta_{M_i}$ is the product of the
$L(s,\chi,M_i/E_i)$ over the characters $\chi$ of the abelian group $H_i/\ker\psi_i$. Each factor with $\chi\neq1$ is
entire, and the factor with $\chi=1$ is $\zeta_{E_i}$, whose only pole is at $s=1$. So every zero of
$L(s,\psi_i)$ with $s\ne1$ is a zero of $\zeta_{M_i}$. By the Aramata--Brauer theorem $\zeta_L/\zeta_{M_i}$ is entire,
so every such zero is also a zero of $\zeta_L$, and hence lies in $\operatorname{Re}s\le\theta$. It follows that in
$\operatorname{Re}s>\theta$ every Brauer factor is holomorphic and nonvanishing away from $s=1$. The order at $s=1$ is
$\langle\rho,1\rangle$ by Frobenius reciprocity.
\end{proof}

\begin{proposition}[$\theta$-Chebotarev]\label{prop:cheb}
There is an absolute effective constant $c_0$ with the following property. Let $L/K$ be a Galois extension of number
fields with group $G$, let $C\subset G$ be a conjugacy class, and assume $\Zt_L(\theta)$. Let $\pi_C(x)$ count
the prime ideals $\fp$ of $K$ that are unramified in $L$, have $\Nm\fp\le x$, and have $\Frob_\fp\in C$. Then for
$x\ge2$,
\[
  \Bigl|\pi_C(x)-\tfrac{|C|}{|G|}\Li(x)\Bigr|
  \le c_0\Bigl[\tfrac{|C|}{|G|}\,x^{\theta}\Bigl(\log(d_Lx^{n_L})+\frac{\log d_L}{(1-\theta)\log x}\Bigr)+\log d_L\,\log\log(3x)\Bigr].
\]
\end{proposition}

For a union of $r$ conjugacy classes one adds the bounds; the result has the same form, with the last term multiplied
by $r$.

\begin{proof}
All implied constants in this proof are absolute and effective; this is the convention of Lagarias--Odlyzko \cite[\S2]{LO},
whose ``$\ll$'' denotes an effectively computable absolute constant.

\emph{The explicit formula.} Fix $g\in C$ and put $E=L^{\langle g\rangle}$. The unconditional explicit formula
\cite[Thm.~7.1, (7.4)--(7.5)]{LO} states that for $x\ge2$ and $T\ge2$,
\begin{multline*}
  \Bigl|\psi_C(x)-\tfrac{|C|}{|G|}x+S(x,T)\Bigr|\ll\tfrac{|C|}{|G|}\Bigl[\tfrac{x\log x+T}{T}\bigl(\log d_L+n_L\log x\bigr)
  +\tfrac{n_Lx\log x\log T}{T}\Bigr]\\+\log x\log d_L+n_K\tfrac{x(\log x)^2}{T},
\end{multline*}
where
\[
  S(x,T)=\frac{|C|}{|G|}\sum_\chi\overline{\chi}(g)\Bigl(\sum_{|\gamma|<T}\frac{x^\rho}{\rho}-\sum_{|\rho|<1/2}\frac1\rho\Bigr),
\]
$\chi$ runs over the characters of $\langle g\rangle$, and $\rho=\beta+i\gamma$ over the nontrivial zeros of the Hecke
$L$-function $L(s,\chi,L/E)$.

\emph{Location of the zeros.} The product of the $L(s,\chi,L/E)$ over $\chi$ is $\zeta_L$, so each $\rho$ is a zero of
$\zeta_L$. Under $\Zt_L(\theta)$ and the functional equation, $1-\theta\le\beta\le\theta$. Hence $|x^\rho|\le x^\theta$
and $|\rho|\ge1-\theta$ for every $\rho$.

\emph{Zeros with $|\rho|<\tfrac12$.} These occur in both inner sums of $S(x,T)$, and are not excluded by
$\Zt_L(\theta)$ when $\theta>\frac12$: they have $1-\theta\le\beta<\frac12$ and $|\gamma|<\frac12$. Each contributes at
most $(x^\theta+1)/(1-\theta)$ in absolute value. Their number is at most $n_\chi(0)$, the number of zeros with
$|\gamma|\le1$, and \cite[Lemma~5.4]{LO} gives $n_\chi(0)\ll\log A(\chi)+n_E$, where $A(\chi)=d_E\Nm\mathfrak f_\chi$.

\emph{Zeros with $|\rho|\ge\tfrac12$ and $|\gamma|<T$.} By \cite[Lemma~5.4]{LO} applied at $t=0,\pm1,\dots,\pm\lfloor T\rfloor$ (a zero counted at $t=j$ has $|\rho|\ge\max(|j|-1,\frac12)$),
\[
  \sum_{\substack{|\gamma|<T\\|\rho|\ge1/2}}\Bigl|\frac{x^\rho}{\rho}\Bigr|\le x^\theta\sum_{|j|\le T}\frac{n_\chi(j)}{\max(|j|-1,\frac12)}
  \ll x^\theta\bigl(\log A(\chi)+n_E\log T\bigr)\log T,
\]
exactly as in the proof of \cite[Thm.~9.1]{LO}, with $x^\theta$ in place of $x^{1/2}$.

\emph{Summing over $\chi$.} The group $\langle g\rangle$ is abelian, so $|\chi(g)|=1$; by the conductor--discriminant
formula $\sum_\chi\log A(\chi)=\log d_L$, and $\sum_\chi n_E=n_L$. Therefore
\[
  |S(x,T)|\ll\tfrac{|C|}{|G|}\Bigl[x^\theta\bigl(\log d_L+n_L\log T\bigr)\log T+\frac{x^\theta(\log d_L+n_L)}{1-\theta}\Bigr].
\]
If $n_L\ge2$, Minkowski's bound gives $\log d_L\ge\frac{n_L}{2}\log3$, so $n_L\ll\log d_L$. If $n_L=1$, then $L=\Q$ and
$\zeta$ has no zeros with $|\rho|<\frac12$ (its first zero has height $>14$), so the last term is absent. In both cases the
last term is $\ll\frac{|C|}{|G|}x^\theta\log d_L/(1-\theta)$.

\emph{Choice of $T$ and passage to $\pi_C$.} Take $T=x$. Since $|C|\ge1$ we have $n_K=n_L/|G|\le\frac{|C|}{|G|}n_L$, so
\[
  \Bigl|\psi_C(x)-\tfrac{|C|}{|G|}x\Bigr|\ll\tfrac{|C|}{|G|}x^\theta\Bigl[\log x\,\log(d_Lx^{n_L})+\frac{\log d_L}{1-\theta}\Bigr]+\log x\log d_L .
\]
As in \cite[\S9]{LO}, replacing $\psi_C$ by $\theta_C(x)=\sum\log\Nm\fp$ (over the unramified $\fp$ counted by
$\pi_C$) changes it by $\ll n_K\sqrt x\log x+\log d_L\log x$ (prime powers, and ramified primes), which is absorbed.
Partial summation, $\pi_C(x)=\theta_C(x)/\log x+\int_2^x\theta_C(t)\,dt/(t\log^2t)$, then gives the claim: the terms
$t^\theta\log(d_Lt^{n_L})\log t$ and $t^\theta\log d_L/(1-\theta)$ integrate to $\ll x^\theta\log(d_Lx^{n_L})$ and
$\ll x^\theta\log d_L/((1-\theta)\log x)$ (using $\theta\ge\frac12$), and the term $\log t\log d_L$ integrates to
$\ll\log d_L\log\log(3x)$.
\end{proof}

\begin{remark}[comparison with the literature]\label{rem:cheb-lit}
Under a $\theta$-hypothesis a statement of the same kind, without explicit $\theta$-dependence, is \cite[Thm.~4]{CM}:
$\pi_C(x,L/K)=\frac{\#C}{[L:K]}\pi(x)+O\bigl(\#C\,x^\theta[K:\Q](\frac{\log d_L}{[L:\Q]}+\log x)\bigr)$, stated with the
remark that Serre's proof for $\theta=\frac12$ \cite[Thm.~4]{Serre81} goes through; the dependence of the implied
constant on $\theta$ is not discussed. David--Wu \cite[Thm.~3.8]{DW} prove, for $D\subset G$ a union of classes with
$HD\subset D$ for a normal subgroup $H$ such that the Artin $L$-functions of $G/H$ satisfy Artin's conjecture and the
$\theta$-hypothesis, $\pi_D(x,L/K)=\frac{|D|}{|G|}\Li(x)+O\bigl((|D|/|H|)^{1/2}x^\theta n_K\log(M(L/K)x)\bigr)$, again
without explicit $\theta$-dependence. In Proposition~\ref{prop:cheb} the dependence on $\theta$ is explicit: the factor
$(1-\theta)^{-1}$ comes only from zeros close to the real axis, and it is negligible once $\log x\gg(1-\theta)^{-1}$.
\cite[p.~3]{CT} gives the case $G=\{1\}$ (the prime ideal theorem). Frei, Sofos and Widmer \cite[Thm.~1.1]{FSW}
prove, assuming only that $\zeta_L$ has no zero in $\operatorname{Re}s>1-\delta$ ($0<\delta\le\frac12$), that
$|\pi_C(x,L/K)-\frac{\#C}{\#G}\Li(x)|\le c_1\frac{\#C}{\#G}x^{1-\delta}\log(\Delta_L^{1/\delta}x^{n_L})$ for $x>2$, with
$c_1$ absolute and effective; with $\theta=1-\delta$ the error is
$\ll\frac{|C|}{|G|}x^\theta\bigl(\frac{\log d_L}{1-\theta}+n_L\log x\bigr)$. This is a statement of the same kind as
Proposition~\ref{prop:cheb}, with a detailed proof (they too note that \cite[Thm.~4]{CM} states a version without one).
In both, the $\theta$-dependence comes from zeros of small modulus (\cite[Lemma~2.3]{FSW}); in Proposition~\ref{prop:cheb}
that term carries an extra factor $1/\log x$. Bhargava \cite[Lemma~7]{Bhv} proves, for $K/\Q$ Galois of degree $n$
with $\Zt_K(1-\delta)$, the upper bound
$\#\{x<p\le2x:p\text{ splits completely in }K\}\ll\frac{x}{n\log x}+\frac{x^{1-\delta}}{\log x}(\frac{\log d_K}{n}+1)$
with an absolute constant, independent of $K$ and $\delta$, and uses it with \cite{CT} to show that every nonsquare
integer $a\ne-1$ is a primitive root modulo a set of primes of positive lower density. The case used in Theorem~A is also proved
independently in Lean (Appendix~\ref{sec:lean}), with a loss of $x^\varepsilon$, via the prime ideal theorems for the two
fields of the quadratic extension.
\end{remark}

%===============================================================================
\section{Theorem A: vanishing of Fourier coefficients}\label{sec:A}
%===============================================================================

\begin{mainthm}\label{thm:A}
Let $f=\sum_{n\ge1}a(n)q^n$ be a normalised cuspidal Hecke newform of weight $k\ge2$ and level $N$, with trivial
character and $a(n)\in\Z$, and without complex multiplication. Trivial character and rational coefficients force $k$
to be even. Let $\tfrac12\le\theta<1$, and assume $\Zt(\theta)$; more precisely, it suffices to assume
$\Zt_M(\theta)$, with $\theta$ fixed, for the fields $M=M_\ell$ defined in the proof, for all sufficiently large
primes $\ell$. Then
\[
  \#\{p\le x:\ a(p)=0\}\ \ll_{f,\theta}\ x^{(1+\theta)/2}.
\]
The implied constant is $\ll_f(1-\theta)^{-1}$. It is effective whenever the finite set $S_f$ of primes with small
residual image (step~1 of the proof) is effectively known, as it is for $f=\Delta$. At $\theta=\frac12$ (GRH) this is
Serre's $x^{3/4}$.
\end{mainthm}

\begin{corollary}[assuming (I2)]\label{cor:A-I2}
Suppose that the Lean theorem \texttt{DedekindZetaZeroFree.main} of \cite{CT-Lean}, displayed in (I2), is correct as a
statement about Dedekind zeta functions, i.e.\ that $\Zt(31/32)$ holds. (The proof of that Lean theorem has been
checked by the Lean kernel; as of 14 October 2026 there is no human-readable proof of the constant $31/32$.) Then, for
$f$ as in Theorem~\ref{thm:A},
\[
  \#\{p\le x:\ a(p)=0\}\ll_f x^{63/64},\qquad\text{in particular}\qquad \#\{p\le x:\ \tau(p)=0\}\ll x^{63/64}.
\]
The bound $x^{63/64}$ is proved here on paper, from Theorem~\ref{thm:A}. What the Lean kernel has checked is the
weaker implication with exponent $(1+\theta)/2+\varepsilon$, i.e.\ $x^{63/64+\varepsilon}$ for every $\varepsilon>0$,
conditional on hypothesis H1 (Appendix~\ref{sec:lean}); the sharp exponent $63/64$ is not formalised.
\end{corollary}

\begin{corollary}[assuming (I1)]\label{cor:A-I1}
If Theorem~1.1 of the unrefereed preprint \cite{CT} holds, i.e.\ $\Zt(1-1/6000)$, then
$\#\{p\le x:\ a(p)=0\}\ll_f x^{1-1/12000}$.
\end{corollary}

\begin{remark}[weight $2$]
For $k=2$, $f$ corresponds to an elliptic curve over $\Q$. Elkies' unconditional bound $x^{3/4}$ \cite{Elk87,Elk91}
for supersingular primes is better than $x^{(1+\theta)/2}$ for every $\theta>\frac12$. For $k\ge4$ the best
unconditional bound we know of is Thorner--Zaman's $x(\log\log x)^2/(\log x)^2$ \cite[Thm.~1.4]{TZ}, a logarithmic
saving over $\pi(x)$.
\end{remark}

\subsection{The group lemma}

Let $\ell\ge5$ be prime and $G'=\PGL_2(\F_\ell)$, acting on $\PP^1(\F_\ell)$ by M\"obius transformations. Let
$T\subset G'$ be the image of the diagonal torus, which is the pointwise stabiliser of $\{0,\infty\}$. Let
$\Nor=N_{G'}(T)$, the setwise stabiliser of $\{0,\infty\}$, and $\Nout=\Nor\setminus T$, the images of
$\left(\begin{smallmatrix}0&\beta\\1&0\end{smallmatrix}\right)$.

\begin{lemma}\label{lem:group}
\begin{enumerate}[label=(\alph*)]
\item $[G':T]=\ell(\ell+1)$, $[G':\Nor]=\ell(\ell+1)/2$ and $[\Nor:T]=2$.
\item If $g\in\GL_2(\F_\ell)$ has $\tr g=0$, then its image $\iota\in G'$ is an involution. Conversely, every
  preimage in $\GL_2(\F_\ell)$ of an involution of $G'$ has trace~$0$.
\item For every involution $\iota\in G'$,
  \[\#\{h\Nor\in G'/\Nor:\ h^{-1}\iota h\in \Nout\}\ \ge\ \tfrac{\ell-1}{2}.\]
  (The count is exactly $(\ell-1)/2$ or $(\ell+1)/2$, according as $\iota$ has two or no fixed points on
  $\PP^1(\F_\ell)$.)
\end{enumerate}
\end{lemma}

\begin{proof}
(a) $G'$ acts simply transitively on ordered triples of distinct points of $\PP^1(\F_\ell)$. So it acts transitively
on ordered pairs of distinct points, and the stabiliser of $(0,\infty)$ is $T$; hence $[G':T]=(\ell+1)\ell$. Similarly
$G'$ acts transitively on unordered pairs, and the stabiliser of $\{0,\infty\}$ is $\Nor$; hence $[G':\Nor]=\ell(\ell+1)/2$.

(b) By Cayley--Hamilton, $g^2=-\det(g)\,I$ is scalar. Moreover $g$ is not scalar, since a scalar matrix $cI$ with
trace $2c=0$ would have $c=0$ ($\ell$ odd). So $\iota\neq1$ and $\iota^2=1$. Conversely, if $g$ is not scalar and $g^2=cI$,
then the minimal polynomial of $g$ is $X^2-c$, which is then its characteristic polynomial, so $\tr g=0$.

(c) The map $h\Nor\mapsto\{h0,h\infty\}$ is a bijection from $G'/\Nor$ to the set of unordered pairs of distinct points.
Now $h^{-1}\iota h\in\Nor$ if and only if $\iota$ preserves $\{h0,h\infty\}$, and $h^{-1}\iota h\in T$ if and only if
$\iota$ fixes both points. So $h^{-1}\iota h\in\Nout$ if and only if $\iota$ swaps $h0$ and $h\infty$. A nontrivial
element of $G'$ has at most two fixed points on $\PP^1$, since a non-scalar matrix has at most two eigenlines. Since
$\iota^2=1$, the remaining $\ge\ell-1$ points split into $2$-cycles of $\iota$. Each $2$-cycle is one swapped pair,
which gives at least $(\ell-1)/2$ of them.
\end{proof}

\subsection{The decomposition dictionary}

Let $L'/\Q$ be Galois with group $G'$ and let $H\le G'$. Let $p$ be unramified in $L'$, fix a prime $\fP'\mid p$ of
$L'$, and let $\sigma$ be its Frobenius. The primes of $L'^H$ above $p$ are the $g\fP'\cap L'^H$, $g\in G'$; they
correspond to double cosets $HgD$, where $D=\langle\sigma\rangle$. The residue degree of $g\fP'\cap L'^H$ over $p$ is
$[gDg^{-1}:H\cap gDg^{-1}]$. Hence:

\begin{lemma}\label{lem:dict}
The primes $\fp$ of $L'^H$ above $p$ with residue degree $1$ correspond bijectively to the right cosets $Hg$ with
$g\sigma g^{-1}\in H$; for these the double coset $HgD$ is the single coset $Hg$. For such a prime, $\Nm\fp=p$, and the
Frobenius of $g\fP'$ over $\fp$ in $\Gal(L'/L'^H)=H$ is $g\sigma g^{-1}$.
\end{lemma}

Writing $h=g^{-1}$, the condition $g\sigma g^{-1}\in H$ becomes $h^{-1}\sigma h\in H$, and $Hg\mapsto hH$ is a
bijection. So in Lemma~\ref{lem:group}(c) we may count either right or left cosets.

\subsection{Proof of Theorem~\ref{thm:A}}

\emph{Step 1 (large image).} By Deligne, for each prime $\ell$ there is a continuous representation
$\bar\rho_\ell=\bar\rho_{f,\ell}:\Gal(\overline\Q/\Q)\to\GL_2(\F_\ell)$ with the following properties: it is unramified
outside $\ell N$; for $p\nmid\ell N$, $\tr\bar\rho_\ell(\Frob_p)\equiv a(p)\pmod\ell$; and
$\det\bar\rho_\ell=\chi_\ell^{k-1}$, where $\chi_\ell$ is the mod-$\ell$ cyclotomic character. Since $f$ has no CM,
Ribet \cite{Rib85} and Momose \cite{Mom81} show that there is a finite set $S_f$ of primes such that the image of
$\bar\rho_\ell$ contains $\SL_2(\F_\ell)$ for all $\ell\notin S_f$. For $f=\Delta$, Swinnerton-Dyer
\cite{SwD73} gives $S_\Delta=\{2,3,5,7,23,691\}$. Since $k-1$ is odd, $\det\bar\rho_\ell$ takes non-square values. An
element of the image with non-square determinant does not map into $\PSL_2(\F_\ell)$. Hence for $\ell\notin S_f$ the
projective image is all of $G'=\PGL_2(\F_\ell)$.

Fix $\ell\ge5$ with $\ell\notin S_f$ and $\ell\nmid N$. Let $L'$ be the field cut out by the projective representation,
so that $\Gal(L'/\Q)=G'$. Put $E=L'^{\Nor}$ and $M=M_\ell=L'^T$. Then $M/E$ is quadratic with group $\Nor/T$, and by
Lemma~\ref{lem:group}(a),
\[
  n_E=\ell(\ell+1)/2,\qquad n_M=\ell(\ell+1).
\]

\emph{Step 2 (reduction to inert primes).} Let $p\nmid\ell N$ with $a(p)=0$. Then $\tr\bar\rho_\ell(\Frob_p)=0$, so
the Frobenius $\sigma$ of $p$ in $G'$ is an involution (Lemma~\ref{lem:group}(b)). By Lemma~\ref{lem:group}(c) and
Lemma~\ref{lem:dict}, at least $(\ell-1)/2$ distinct primes $\fp$ of $E$ lie above $p$ with $\Nm\fp=p$ and with
Frobenius in $\Gal(L'/E)=\Nor$ lying in $\Nout$. Such a $\fp$ is unramified in $M$, and since $\Nout$ maps to the
nontrivial element of $\Nor/T=\Gal(M/E)$, it is inert in $M$. Let $\pi_\inert(x)$ be the number of prime ideals of $E$
that are unramified and inert in $M$ and have norm $\le x$. Counting pairs $(p,\fp)$ gives
\begin{equation}\label{eq:reduction}
  \#\{p\le x:\ a(p)=0\}\ \le\ \omega(\ell N)+\frac{2}{\ell-1}\,\pi_\inert(x).
\end{equation}

\emph{Step 3 (Chebotarev in one quadratic extension).} We apply Proposition~\ref{prop:cheb} to $M/E$, with
$G=\Z/2$ and $C$ the nontrivial element, so $|C|/|G|=\frac12$. This needs only $\Zt_M(\theta)$, and it gives
\[
  \pi_\inert(x)\le\tfrac12\Li(x)+c_0\Bigl[\tfrac12x^\theta\Bigl(\log d_M+n_M\log x+\frac{\log d_M}{(1-\theta)\log x}\Bigr)+\log d_M\log\log(3x)\Bigr].
\]
The field $M$ has degree $n_M=\ell(\ell+1)$ and is unramified outside $\ell N$. Serre's bound \cite[Prop.~6]{Serre81},
which holds for every number field, gives
\[
  \log d_M\le(n_M-1)\sum_{p\mid\ell N}\log p+n_M\,\omega(\ell N)\log n_M\le\ell(\ell+1)\bigl(\log(\ell N)+3\,\omega(\ell N)\log\ell\bigr)\ll_N\ell^2\log\ell .
\]
Inserting this, for $x\ge e$,
\[
  \pi_\inert(x)\le\tfrac12\Li(x)+\frac{c_1(f)}{1-\theta}\,\bigl[\ell^2x^\theta\log(\ell x)+\ell^2\log\ell\,\log\log(3x)\bigr].
\]

\emph{Step 4 (choice of $\ell$).} Combining this with \eqref{eq:reduction},
\[
  \#\{p\le x:\ a(p)=0\}\le\frac{\Li(x)}{\ell-1}+\frac{c_2(f)}{1-\theta}\bigl[\ell x^\theta\log(\ell x)+\ell\log\ell\,\log\log(3x)\bigr]+\omega(\ell N).
\]
Let $y=x^{(1-\theta)/2}/\log x$. For $x\ge x_0(f,\theta)$, Bertrand's postulate gives a prime $\ell\in(y,2y]$, and for
$x$ large this $\ell$ is larger than every element of $S_f$ and every prime divisor of $N$. With this choice,
$\Li(x)/(\ell-1)\ll x/(y\log x)=x^{(1+\theta)/2}$ and $\ell x^\theta\log(\ell x)\ll yx^\theta\log x=x^{(1+\theta)/2}$.
The remaining terms are smaller. \qed

\begin{remark}[why uniformity is essential]\label{rem:uniformity}
The prime $\ell$ grows like $x^{(1-\theta)/2}$, so the fields $M_\ell$ change with $x$. Their degrees are
$\ell(\ell+1)$ and their discriminants are unbounded. A zero-free half-plane for each field separately, with $\theta$
depending on the field, would give nothing. What is needed is exactly the uniformity in $K$ asserted by (I1) and (I2);
classical zero-free regions do not have it. We checked that the uniformity of (I1) and (I2) is genuine as stated:
there is no height cutoff and no exceptional field.
\end{remark}

\begin{remark}[logarithmic factors]\label{rem:smoothing}
Zywina's smoothed Chebotarev theorem \cite[Thms.~2.2--2.3]{Zyw15} gives $x^{3/4}(\log x)^{-1/2}$ under GRH. We expect,
but have not checked, that it transfers to $\Zt(\theta)$ and saves a factor $(\log x)^{1/2}$ in Theorem~\ref{thm:A};
this would not change the exponent, and nothing below depends on it.
\end{remark}

For $\tau$, Theorem~\ref{thm:A} bears only on the density side of Lehmer's question \cite{Leh47}: it shows that the
primes with $\tau(p)=0$, if any, have density zero with a power saving.

%===============================================================================
\section{Theorem B: an uncertainty principle for involutions, and three error functionals}\label{sec:B}
%===============================================================================

Fix $f$ as in Theorem~\ref{thm:A}. For a prime $\ell$ write $G_\ell=\PGL_2(\F_\ell)$, and let
$\Inv_\ell=D_s\sqcup D_n$ be the union of the split and nonsplit involution classes. These have sizes
$|D_s|=\ell(\ell+1)/2$ and $|D_n|=\ell(\ell-1)/2$. So $|\Inv_\ell|=\ell^2$, and we put
$\mu_0:=|\Inv_\ell|/|G_\ell|=\ell/(\ell^2-1)$. We write $T_s$ and $T_n$ for a split and a nonsplit maximal torus of
$G_\ell$ (cyclic of orders $\ell-1$ and $\ell+1$), $\Nor(T_s)$ and $\Nor(T_n)$ for their normalisers (dihedral of
orders $2(\ell-1)$ and $2(\ell+1)$), and $B$ for a Borel subgroup (of order $\ell(\ell-1)$).

For a real class function $\varphi$ on a finite group $G$, let $\mu(\varphi)=|G|^{-1}\sum_g\varphi(g)$ and
$\widehat\varphi(\pi)=|G|^{-1}\sum_g\varphi(g)\overline{\chi_\pi(g)}$. Bella\"iche's \emph{Littlewood norm}
\cite{Bel16} is
\[
  \lambda_A(\varphi)=\sum_{\pi\in\widehat G}|\widehat\varphi(\pi)|\dim\pi,
\]
which is the Fourier-algebra norm of $\varphi$. Write $\lambda_{\mathrm{mon}}(\varphi)$ for the least value of
$\sum|c_i|[G:H_i]$ over decompositions $\varphi=\sum c_i\Ind_{H_i}^G\psi_i$ into induced linear characters; then
$\lambda_{\mathrm{mon}}\ge\lambda_A$. A \emph{majorant} of $D\subset G$ is a function $\varphi\ge0$ with $\varphi\ge1$
on $D$.

\subsection{The three classes}\label{subsec:B-classes}

If $p\nmid\ell N$ and $a(p)=0$, then the projective Frobenius of $p$ lies in $\Inv_\ell$
(Lemma~\ref{lem:group}(b)). The arguments we consider use only this information, through one or several primes
$\ell$. They must therefore be valid at \emph{every} element of $\Inv_\ell$ (respectively of a product of such sets),
and not only at the Frobenius classes of the primes with $a(p)=0$, about which nothing more is known
unconditionally. This requirement is part of the following definition. (Without it the definition would be vacuous:
if $\{p:a(p)=0\}$ were finite, as is conjectured for $k\ge4$, every count would be a valid certificate.)

For a finite set $P$ of primes $\ell\notin S_f$ with $\ell\nmid N$, let $L'_\ell$ be the field cut out by the
projective representation attached to $\bar\rho_{f,\ell}$, put $\Gamma_P=\prod_{\ell\in P}G_\ell$ and
$S_P=\prod_{\ell\in P}\Inv_\ell\subset\Gamma_P$, and say that $P$ is \emph{independent} if the projective image of
$\bigoplus_{\ell\in P}\bar\rho_{f,\ell}$ is all of $\Gamma_P$. (By Ribet \cite{Rib85} and Loeffler \cite{Loe17}, this
holds whenever $P$ avoids a finite set depending on $f$.)

\begin{definition}\label{def:classes}
Fix $\frac12\le\theta<1$, a constant $c>0$, a function $w>0$ with $w(x)\ge x^{-o(1)}$ as $x\to\infty$ (the error
weight; for instance $w\ge1$, or the weight $w\asymp\log\ell/\log x$ of \S\ref{subsec:B-scope}), and a norm
$\lambda\in\{\lambda_A,\lambda_{\mathrm{mon}}\}$.
\begin{enumerate}[label=(\alph*),leftmargin=*]
\item \emph{Single-$\ell$ majorants.} The data are a prime $\ell=\ell(x)\ge5$ with $\ell\notin S_f$, and a class
  function $\varphi$ on $G_\ell$ that is a majorant of $\Inv_\ell$. The certified bound is
  \[
    B_{\mathrm{maj}}(x)=\mu(\varphi)\Li(x)+c\,\lambda(\varphi)\,x^\theta\,w(x).
  \]
\item \emph{Single-extension certificates.} The data are: an independent set $P$ of primes $\ell\ge7$; a finite
  Galois extension $\Omega/\Q$ containing $L'_\ell$ for every $\ell\in P$, with group $\cG=\Gal(\Omega/\Q)$ and
  projection $\pi:\cG\to\Gamma_P$; fields $K\subset L\subset\Omega$ with $L/K$ Galois, $H=\Gal(\Omega/K)$ and
  $N_0=\Gal(\Omega/L)\lhd H$; a conjugation-stable $C\subset\Gal(L/K)=H/N_0$ with preimage $\widetilde C\subset H$; and
  an integer $m\ge1$. Put $\widetilde S=\pi^{-1}(S_P)$, the full preimage of $\prod_\ell\Inv_\ell$ (if $\Omega$
  contains the fields cut out by the $\bar\rho_{f,\ell}$, then by Lemma~\ref{lem:group}(b) this is the set of $g$ with
  $\tr\bar\rho_{f,\ell}(g)=0$ for all $\ell\in P$). The data must satisfy the \emph{validity condition}
  \[
    \#\{xH\in\cG/H:\ x^{-1}gx\in\widetilde C\}\ \ge\ m\qquad\text{for every }g\in\widetilde S.
  \]
  The certified bound is
  \[
    B_{\mathrm{cert}}(x)=\cA\,\Li(x)+c\,\cE\,x^\theta\log x,\qquad
    \cA=\frac{|C|}{|\Gal(L/K)|\,m},\quad \cE=\frac{|C|\,n_K}{m}.
  \]
\item \emph{Per-prime Selberg sieves.} The data are: a real $\eta>0$ and $\gamma\ge1$; an independent set $L^*$ of
  primes $5\le\ell\le x^\gamma$; for each $\ell\in L^*$ a real class function $v_\ell$ on $G_\ell$ with
  \[
    \mathbb E v_\ell=0,\qquad \mathbb E v_\ell^2=1,\qquad \tilde r_\ell:=(\mathbb E_{\Inv}v_\ell)^2\ge1+\eta,
  \]
  where $\mathbb E$ is the mean over $G_\ell$ and $\mathbb E_{\Inv}$ the mean over $\Inv_\ell$; and a finitely
  supported real vector $y=(y_m)$, indexed by squarefree $m$ supported on $L^*$, such that
  $P:=\sum_my_mv_m\ge1$ on $S_{L^*}$, where $v_m=\bigotimes_{\ell\mid m}v_\ell$ ($v_1=1$). The certified bound is
  \[
    B_{\mathrm{sieve}}(y)=\Li(x)\sum_my_m^2+c\,x^\theta w(x)\sum_{m,n}|y_my_n|\,E(m,n),\qquad E(m,n)=\lambda(v_mv_n)-\delta_{mn}.
  \]
\end{enumerate}
\end{definition}

\begin{remark}[what the functionals represent]\label{rem:soundness}
Each functional is the bound that the corresponding argument proves, with its error term in a fixed, postulated form.
\begin{enumerate}[label=(\alph*),leftmargin=*]
\item Under $\Zt(\theta)$, for $\varphi=\Ind_H^G\psi$ with $\psi$ linear, $\sum_p\varphi(\Frob_p)$ is, up to
  $O([G:H]x^{1/2})$, a sum of $\psi(\Frob_\fp)$ over degree-one primes $\fp$ of $L^H$. The Hecke $L$-function
  $L(s,\psi,L/L^H)$ divides $\zeta_L$ (Aramata--Brauer), so the argument of Proposition~\ref{prop:cheb} applies to it,
  with $\log(d_{L^H}\Nm\mathfrak f_\psi)\ll[G:H]\log(q|G|)$ by \cite[Prop.~6]{Serre81}, if $L$ is unramified outside
  the primes dividing $q$. By linearity, for every class function $\varphi$ on $G=\Gal(L/\Q)$,
  \begin{multline}\label{eq:majorant-form}
    \Bigl|\sum_{p\le x,\,p\nmid q}\varphi(\Frob_p)-\mu(\varphi)\Li(x)\Bigr|\ll\lambda_{\mathrm{mon}}(\varphi)\,x^\theta\log(q|G|x)\Bigl(1+\frac1{(1-\theta)\log x}\Bigr)\\
    +\lambda_{\mathrm{mon}}(\varphi)\log(q|G|)\log\log(3x).
  \end{multline}
  Under GRH together with Artin's conjecture, Bella\"iche \cite[Thm.~1]{Bel16} proves the analogue with $\lambda_A$
  and $x^{1/2}$. Since $\varphi\ge0$, and $\varphi\ge1$ on all of $\Inv_\ell$, $\#\{p\le x:a(p)=0\}\le\sum_{p\le x,\,p\nmid q}\varphi(\Frob_p)+\omega(q)$,
  and \eqref{eq:majorant-form} bounds this by $B_{\mathrm{maj}}$ plus lower-order terms, with $\lambda=\lambda_{\mathrm{mon}}$
  and $w(x)=\log(q|G|x)$ (a sharp cutoff); smoothing changes $w$ (\S\ref{subsec:B-scope}).
\item By Lemma~\ref{lem:dict}, every $p$ unramified in $\Omega$ with $a(p)=0$ has its Frobenius class in $\widetilde
  S$, hence, by validity, at least $m$ primes $\fp\mid p$ of $K$ of degree one with $\Frob_\fp\in C$.
  Proposition~\ref{prop:cheb} applied to $L/K$ (summed over the classes in $C$) then gives $\#\{p\le x:a(p)=0\}\le\cA\Li(x)+c_0\cE x^\theta\log x$ plus
  further error terms, which only increase the bound. We call $\cA\Li(x)+c\,\cE x^\theta\log x$ the
  \emph{Lagarias--Odlyzko (LO) error shape}. The certificate of \S\ref{sec:A} has $K=E_\ell$, $|C|=1$ and
  $m=(\ell-1)/2$.
\item $P^2$ is a majorant of $S_{L^*}$, and the $v_m$ are orthonormal, so $\sum_{p\le x}P^2(\Frob_p)$ has main term
  $\Li(x)\sum_my_m^2$. Bounding the error term of each product $v_mv_n$ separately, by the analogue of
  \eqref{eq:majorant-form}, gives $B_{\mathrm{sieve}}$. With $\lambda=\lambda_{\mathrm{mon}}$ this needs only $\Zt(\theta)$
  for the composita; with $\lambda=\lambda_A$ it also needs Artin's conjecture. The functional postulates
  \emph{pairwise} error accounting: no cancellation between the error terms of different pairs is allowed.
\end{enumerate}
Since $\lambda_{\mathrm{mon}}\ge\lambda_A$, each functional is smallest for $\lambda=\lambda_A$; a lower bound proved
for $\lambda=\lambda_A$ therefore holds for both choices.
\end{remark}

\subsection{An exact Littlewood-norm inequality}

\begin{thmB}[exact uncertainty for involutions]\label{thm:B1}
Let $\ell\ge5$ be prime, $G=\PGL_2(\F_\ell)$, and let $\varphi$ be a class-function majorant of $\Inv_\ell$. Then
\[
  \lambda_A(\varphi)\ \ge\ \max\Bigl\{\lambda_0-2(\ell-1)\bigl(\mu(\varphi)-\mu_0\bigr),\ \ \ell-(\ell-1)\mu(\varphi),\ \ \mu(\varphi)\Bigr\},\qquad
  \lambda_0:=\lambda_A(1_{\Inv})=\frac{\ell^2+1}{\ell+1}.
\]
Consequently
\[
  \mu(\varphi)\,\lambda_A(\varphi)\ \ge\ \kappa_\ell:=\mu_0\lambda_0=\frac{\ell(\ell^2+1)}{(\ell+1)(\ell^2-1)}\ \ge\ 1-\frac1\ell,
\]
with equality for $\varphi=1_{\Inv}$; and $\mu(\varphi)\lambda_A(\varphi)\ge1$ if $\mu(\varphi)\ge\frac12$.
\end{thmB}

\begin{proof}
\emph{Characters.} We use the standard character table of $\GL_2(\F_\ell)$ (see \cite[\S5.2]{FH} or \cite{PS83});
the irreducible characters of $G$ are those of $\GL_2(\F_\ell)$ that are trivial on the centre. They are: $1$; $\mathrm{sgn}$ (the Legendre symbol of the
determinant); the Steinberg character $\mathrm{St}$ (dimension $\ell$) and $\mathrm{sgn}\cdot\mathrm{St}$; the
principal series $\chi_\xi$ (dimension $\ell+1$), indexed by the $(\ell-3)/2$ pairs $\{\xi,\xi^{-1}\}$ of characters of
$\F_\ell^\times$ with $\xi^2\ne1$; and the discrete series $\pi_\psi$ (dimension $\ell-1$), indexed by the $(\ell-1)/2$
pairs $\{\psi,\psi^{-1}\}$ of characters of the nonsplit torus $T_n\cong\Z/(\ell+1)$ with $\psi^2\ne1$. All are
real-valued. Let $\varsigma=\bigl(\frac{-1}{\ell}\bigr)$, so that $\mathrm{sgn}(\iota_s)=\varsigma$ and
$\mathrm{sgn}(\iota_n)=-\varsigma$ for $\iota_s\in D_s$, $\iota_n\in D_n$.

\emph{Fourier coefficients of $1_{\Inv}$.} A direct computation shows that every coefficient is nonzero. Their signs
and absolute values are:
\begin{itemize}
\item $1$: sign $+$, absolute value $\mu_0$; $\ \mathrm{sgn}$: sign $\varsigma$, absolute value $1/(\ell^2-1)$;
\item $\mathrm{St}$: sign $+$, absolute value $1/(\ell^2-1)$; $\ \mathrm{sgn}\cdot\mathrm{St}$: sign $\varsigma$, absolute value $\mu_0$;
\item $\chi_\xi$: sign $\xi(-1)$, absolute value $1/(\ell-1)$; $\ \pi_\psi$: sign $-\psi(\iota_n)$, absolute value $1/(\ell+1)$.
\end{itemize}
Hence
\[
\lambda_0=\frac{\ell+1}{\ell-1}+\frac{\ell-3}{2}\cdot\frac{\ell+1}{\ell-1}+\frac{\ell-1}{2}\cdot\frac{\ell-1}{\ell+1}
=\frac{\ell+1}{2}+\frac{(\ell-1)^2}{2(\ell+1)}=\frac{\ell^2+1}{\ell+1}.
\]
In particular $|\widehat{1_{\Inv}}(\pi)|\le\dim\pi/(\ell^2-1)$ for every $\pi\ne1$, with equality for $\mathrm{sgn}$,
$\mathrm{sgn}\cdot\mathrm{St}$, $\chi_\xi$ and $\pi_\psi$.

\emph{Dual functions.} If $\Psi=\sum_\pi c_\pi\dim\pi\,\chi_\pi$ with real $|c_\pi|\le1$, then for every real class
function $\varphi$
\[
  \lambda_A(\varphi)\ \ge\ \sum_\pi c_\pi\widehat\varphi(\pi)\dim\pi=\mathbb{E}_G[\varphi\Psi].
\]
We use three such functions.
(1)~$\Psi=\sum_\pi\varepsilon_\pi\dim\pi\,\chi_\pi$, where $\varepsilon_\pi$ is the sign of $\widehat{1_{\Inv}}(\pi)$.
Evaluating $\Psi$ from the character table gives $(\ell-1)^2$ at the identity; $-(\ell-1)$ at a nontrivial unipotent
element; $\ell^2-1$ at a split involution; $(\ell-1)^2$ at a nonsplit involution; $0$ at every other split regular
element; and $-2(\ell-1)$ at every other nonsplit regular element. In particular $\Psi\ge0$ on $\Inv$ and $\Psi\ge-2(\ell-1)$ everywhere,
and $\mathbb{E}[1_{\Inv}\Psi]=(|D_s|(\ell^2-1)+|D_n|(\ell-1)^2)/|G|=\lambda_0$.
(2)~$\Psi'=(\ell^2-1)1_{\Inv}-(\ell-1)$. Then $\widehat{\Psi'}(1)=(\ell^2-1)\mu_0-(\ell-1)=1$ and
$|\widehat{\Psi'}(\pi)|=(\ell^2-1)|\widehat{1_{\Inv}}(\pi)|\le\dim\pi$ for $\pi\ne1$.
(3)~$\Psi''=\chi_1=1$.

\emph{The three bounds.} Since $\varphi-1_{\Inv}\ge0$ and $\Psi\ge-2(\ell-1)$,
\[
  \lambda_A(\varphi)\ge\mathbb{E}[1_{\Inv}\Psi]+\mathbb{E}[(\varphi-1_{\Inv})\Psi]\ge\lambda_0-2(\ell-1)(\mu(\varphi)-\mu_0).
\]
Since $\varphi\ge1$ on $\Inv$, $\mathbb E[\varphi1_{\Inv}]\ge\mu_0$, so
$\lambda_A(\varphi)\ge\mathbb E[\varphi\Psi']\ge(\ell^2-1)\mu_0-(\ell-1)\mu(\varphi)=\ell-(\ell-1)\mu(\varphi)$.
Finally $\lambda_A(\varphi)\ge\mathbb E[\varphi]=\mu(\varphi)$.

\emph{The product.} Write $\mu=\mu(\varphi)$; note $\mu\ge\mu_0$ because $\varphi\ge1_{\Inv}$. We have
$\lambda_0+2(\ell-1)\mu_0=\ell+1$, so the first bound gives $\mu\lambda_A\ge F(\mu):=\mu(\ell+1-2(\ell-1)\mu)$. The
function $F$ is concave, with $F(\mu_0)=\kappa_\ell$ and $F(\tfrac12)=1>\kappa_\ell$; hence $\mu\lambda_A\ge\kappa_\ell$
for $\mu\in[\mu_0,\frac12]$. For $\mu\in[\frac12,1]$ the second bound gives $\mu\lambda_A\ge\mu(\ell-(\ell-1)\mu)$, a
concave function of $\mu$ with values $(\ell+1)/4>1$ at $\frac12$ and $1$ at $1$, hence $\ge1$. For $\mu\ge1$ the third
bound gives $\mu\lambda_A\ge\mu^2\ge1$. Finally, $\kappa_\ell<1$ is equivalent to $\ell^2-2\ell-1>0$, and
$\kappa_\ell\ge1-1/\ell$ is equivalent to $\ell^2(\ell^2+1)\ge(\ell^2-1)^2$.
\end{proof}

\begin{remark}[sharpness, equality, and numerical confirmation]\label{rem:B1-num}
\emph{The three-term bound is sharp.} For $\mu\ge\mu_0$ let $\lambda^*(\mu)$ be the least value of $\lambda_A$ over
class-function majorants of $\Inv$ with mean $\mu$. Then $\lambda^*(\mu)$ equals the three-term bound of
Theorem~\ref{thm:B1} for every $\mu\ge\mu_0$. That bound is piecewise linear: it equals the first branch on
$[\mu_0,\mu_1]$, where $\mu_1:=1/(\ell-1)$, the second on $[\mu_1,1]$ and the third on $[1,\infty)$, with values
$\lambda_0$, $\ell-1$ and $1$ at $\mu_0$, $\mu_1$ and $1$ (the first two branches meet at $\mu_1$ because
$\lambda_0+2(\ell-1)\mu_0=\ell+1$). Since $\lambda_A$ is a norm and the mean is linear, if majorants $\varphi_a$ and
$\varphi_b$ attain the bound at the two ends of one of these intervals, then $t\varphi_a+(1-t)\varphi_b$ is a majorant
with $\lambda_A\le t\lambda_A(\varphi_a)+(1-t)\lambda_A(\varphi_b)$, which is the (linear) bound at its mean; together
with Theorem~\ref{thm:B1} this gives equality on the whole interval. So it suffices to exhibit majorants attaining the
bound at $\mu_0$, at $\mu_1$, and at every $\mu\ge1$. At $\mu_0$ take $1_{\Inv}$, and for $\mu\ge1$ the constant
$\varphi\equiv\mu$, which has $\lambda_A=\mu$. At $\mu_1$ take
\[
  \varphi_1=1_{\Inv}+\frac{2}{(\ell-1)^2}\,1_{N},
\]
where $N$ is the set of nonsplit regular elements that are not involutions. Each of the $\ell(\ell-1)/2$ nonsplit tori
contains $\ell-1$ such elements, so $|N|=\ell(\ell-1)^2/2$ and $\mu(\varphi_1)=\mu_0+1/(\ell^2-1)=\mu_1$. At an element of
$N$ the characters take the values $1$; $\mathrm{sgn}$; $-1$ (St); $-\mathrm{sgn}$ ($\mathrm{sgn}\cdot\mathrm{St}$);
$0$ ($\chi_\xi$); and $-\psi(u)-\psi(u)^{-1}$ ($\pi_\psi$, at the element $u$ of $T_n$). Summing over $N$ by orthogonality
on the cyclic group $T_n$ gives the Fourier coefficients of $\varphi_1$: the coefficient at St vanishes, those at
$\chi_\xi$ are those of $1_{\Inv}$, and every other coefficient has the sign $\varepsilon_\pi$ of the corresponding
coefficient of $1_{\Inv}$, with absolute value
\[
  \begin{array}{l|cccc}
    \pi & 1 & \mathrm{sgn} & \mathrm{sgn}\cdot\mathrm{St} & \pi_\psi\\[2pt] \hline\rule{0pt}{16pt}
    \varsigma\text{, resp.\ }\psi(\iota_n)=+1 & \dfrac1{\ell-1} & \dfrac1{\ell^2-1} & \dfrac{\ell}{\ell^2-1} & \dfrac{\ell-3}{(\ell-1)^2}\\[8pt]
    \varsigma\text{, resp.\ }\psi(\iota_n)=-1 & \dfrac1{\ell-1} & \dfrac1{(\ell-1)^2} & \dfrac{\ell-2}{(\ell-1)^2} & \dfrac1{\ell+1}
  \end{array}
\]
(for sgn and $\mathrm{sgn}\cdot\mathrm{St}$ the row is given by $\varsigma$, for $\pi_\psi$ by $\psi(\iota_n)$); all are
positive for $\ell\ge5$. Hence $\lambda_A(\varphi_1)=\sum_\pi\varepsilon_\pi\widehat{\varphi_1}(\pi)\dim\pi=\mathbb E[\varphi_1\Psi]$,
and since $\Psi=-2(\ell-1)$ on $N$, this is $\lambda_0-2(\ell-1)(\mu_1-\mu_0)=\ell-1$, the value of the bound at $\mu_1$.
The minimisers are therefore explicit: $t1_{\Inv}+(1-t)\varphi_1$ on $[\mu_0,\mu_1]$, $t\varphi_1+(1-t)$ on $[\mu_1,1]$,
and constants beyond. (Computer checks of these coefficients, symbolic for general $\ell$, and of
$\lambda_A(\varphi_1)=\ell-1$ by a computation that does not use the character table, are listed in
\S\ref{sec:computations}.)

\emph{Equality.} Equality $\mu(\varphi)\lambda_A(\varphi)=\kappa_\ell$ holds only for $\varphi=1_{\Inv}$: the concave
function $F$ of the proof lies strictly above $\kappa_\ell$ on $(\mu_0,\frac12]$, and $\mu\lambda_A\ge1>\kappa_\ell$ for
$\mu\ge\frac12$; so equality forces $\mu(\varphi)=\mu_0$, and then $\varphi\ge1_{\Inv}$ forces $\varphi=1_{\Inv}$.

\emph{Numerical confirmation.} We verified the character values, the coefficients of $1_{\Inv}$ and the values of
$\Psi$ numerically for $\ell\le97$, with two independently written implementations (\S\ref{sec:computations}). Linear
programs, solved in floating point with HiGHS with the mean fixed exactly, give for every prime $\ell\le97$ the value
$\lambda^*(\mu)$ on grids covering $\mu_0\le\mu\le3$; at every grid point it agrees with the three-term bound to within
$10^{-9}$, as the argument above predicts (for example $\kappa_5=0.9028$ and $\kappa_{97}=0.9900$). These computations
are confirmation only; nothing above rests on them.
Since $\lambda_{\mathrm{mon}}\ge\lambda_A$, Theorem~\ref{thm:B1} holds for $\lambda_{\mathrm{mon}}$ as well.
\end{remark}

\subsection{Single-$\ell$ majorants}

\begin{corB}[lower bound for $B_{\mathrm{maj}}$]\label{cor:B2}
For every choice of data in Definition~\ref{def:classes}(a), where $\ell=\ell(x)$ may depend on $x$,
\[
  B_{\mathrm{maj}}(x)\ \ge\ 2\sqrt{c\,\kappa_\ell\,\Li(x)\,x^\theta\,w(x)} .
\]
In particular, since $w(x)\ge x^{-o(1)}$, for each fixed $\delta>0$ no choice of $\ell$ and $\varphi$ gives
$B_{\mathrm{maj}}(x)\le x^{(1+\theta)/2-\delta}$ for all large $x$. The same holds if $\varphi$ is a class function on
$\Gal(\Omega/\Q)$, for any finite Galois $\Omega/\Q$ containing the field cut out by $\bar\rho_{f,\ell}$, and is a
majorant of the full preimage of $\{\tr=0\}$.
\end{corB}

\begin{proof}
By Theorem~\ref{thm:B1}, $\lambda(\varphi)\ge\lambda_A(\varphi)\ge\kappa_\ell/\mu(\varphi)$, and by the AM--GM
inequality $B_{\mathrm{maj}}\ge\mu\Li(x)+c\kappa_\ell x^\theta w(x)/\mu\ge2\sqrt{c\kappa_\ell\Li(x)x^\theta w(x)}$.
Since $\kappa_\ell\ge0.9$ for $\ell\ge5$ ($\kappa_5=0.9028$, $\kappa_7=0.9115$, and $\kappa_\ell\ge1-1/\ell\ge0.9$
for $\ell\ge11$), the right-hand side is $x^{(1+\theta)/2-o(1)}$. For the last assertion, convolution with the uniform
probability measure on a normal subgroup contracts every Fourier coefficient, so it does not increase $\lambda_A$; it
fixes $\mu$; and it preserves the majorant property when the bad set is a union of cosets of that subgroup. Averaging
over the kernel of $\Gal(\Omega/\Q)\to G_\ell$, and over the scalars of $\GL_2(\F_\ell)$, therefore reduces to $G_\ell$; and since
$k-1$ is odd, every projective involution has a trace-zero lift in the image (\S\ref{sec:A}, Step~1).
\end{proof}

The certificate of \S\ref{sec:A} corresponds to $\varphi_A=1_{D_s}+\frac{\ell+1}{\ell-1}1_{D_n}$, which has
$\mu\lambda_A=1+1/(\ell-1)$. So it is optimal up to a factor $1+O(1/\ell)$ \emph{as a finite-group quantity}. This
says nothing about logarithmic factors in the final bound (Remark~\ref{rem:smoothing}).

\subsection{Single-extension certificates}

\begin{thmB}[lower bound for $B_{\mathrm{cert}}$]\label{thm:B3}
Consider data as in Definition~\ref{def:classes}(b), and assume in addition that the kernel of
$\pi:\cG\to\Gamma_P$ is central in $\cG$. This holds, for example, whenever $\Omega$ is contained in the compositum of
the fields cut out by the $\bar\rho_{f,\ell}$, $\ell\in P$, since that kernel is then a group of scalar matrices. Then
\[
  B_{\mathrm{cert}}(x)=\cA\,\Li(x)+c\,\cE\,x^\theta\log x\ \ge\ c_1\,(\log\log x)^{-1/2}\,x^{(1+\theta)/2}\qquad(x\ge x_0),
\]
where $c_1>0$ depends only on $c$, and $x_0$ only on $c$ and $\theta$. In particular no such certificate gives an exponent below $(1+\theta)/2$.
\end{thmB}

The proof rests on four lemmas.

\begin{lemma}[coset concentration]\label{lem:conc}
Let $\cG$ be a finite group, $H\le\cG$, $N_0\lhd H$, $\widetilde C\subset H$ a union of $N_0$-cosets stable under
conjugation by $H$, $\Sigma\subset\cG$ stable under conjugation, and $m\ge1$. Suppose that every $g\in\Sigma$ has at
least $m$ cosets $xH$ with $x^{-1}gx\in\widetilde C$. Put $\cA=|\widetilde C|/(|H|m)$,
$\cE=|\widetilde C|[\cG:H]/(|N_0|m)$ and $\nu=\max_y|yN_0\cap\Sigma|$. Then
\[
  \cA\ge\frac{|\Sigma||\widetilde C|}{|\cG||\widetilde C\cap\Sigma|},\qquad \cE\ge\frac{|\Sigma|}{|N_0|},\qquad \cA\cE\ge\frac{|N_0||\Sigma|^2}{\nu^2|\cG|}.
\]
\end{lemma}

In the situation of Definition~\ref{def:classes}(b), $\cA$ and $\cE$ as defined here agree with those of the
definition, since $|C|=|\widetilde C|/|N_0|$, $|\Gal(L/K)|=|H|/|N_0|$ and $n_K=[\cG:H]$.

\begin{proof}
Summing the validity condition over $g\in\Sigma$ gives
$m|\Sigma|\le\sum_{xH}|x\widetilde Cx^{-1}\cap\Sigma|=[\cG:H]|\widetilde C\cap\Sigma|$. Substituting this into the
formulas for $\cA$ and $\cE$ gives the first bound, and also $\cE\ge|\widetilde C||\Sigma|/(|N_0||\widetilde C\cap
\Sigma|)\ge|\Sigma|/|N_0|$. Since $\widetilde C$ is a union of $N_0$-cosets, each meeting $\Sigma$ in at most $\nu$
points, $|\widetilde C|/|\widetilde C\cap\Sigma|\ge|N_0|/\nu$. Multiplying the two lower bounds gives the third.
\end{proof}

\begin{lemma}[reductions]\label{lem:red}
\begin{enumerate}[label=(\roman*)]
\item Let $\varpi:\cG\to\Gamma$ be a surjective homomorphism of finite groups, $S\subset\Gamma$ stable under
  conjugation, $\Sigma=\varpi^{-1}(S)$, and $\bar N=\varpi(N_0)$. Under the hypotheses of Lemma~\ref{lem:conc},
  \[
    \cE\ge\frac{|S|}{|\bar N|},\qquad \cA\cE\ge\frac{|\bar N||S|^2}{\bar\nu^2|\Gamma|},\qquad\bar\nu=\max_{y\in\Gamma}|y\bar N\cap S|.
  \]
\item If $\mathcal P\lhd\cG$ and $\mathcal P\subset N_0$, then validity on $\Sigma$ implies validity on $\Sigma\mathcal P$.
\item (Goursat; $\ell\ge7$.) If $\bar N\le\Gamma_P$ and $\pi_\ell(\bar N)\supseteq\PSL_2(\F_\ell)$, where $\pi_\ell$ is
  the projection to the coordinate $\ell$, then $\bar N$ contains $\PSL_2(\F_\ell)$, placed in coordinate $\ell$.
\end{enumerate}
\end{lemma}

\begin{proof}
(i) Let $K_\varpi=\ker\varpi$. Then $|\Sigma|=|K_\varpi||S|$ and $|N_0|=|\bar N||N_0\cap K_\varpi|$. A coset $yN_0$ maps onto
$\varpi(y)\bar N$ with fibres of size $|N_0\cap K_\varpi|$, and $yN_0\cap\Sigma$ lies over $\varpi(y)\bar N\cap S$; hence
$\max_y|yN_0\cap\Sigma|\le|N_0\cap K_\varpi|\,\bar\nu$. Lemma~\ref{lem:conc} gives
$\cE\ge|K_\varpi||S|/(|\bar N||N_0\cap K_\varpi|)\ge|S|/|\bar N|$ and
$\cA\cE\ge|\bar N||N_0\cap K_\varpi||K_\varpi|^2|S|^2/(|N_0\cap K_\varpi|^2\bar\nu^2|K_\varpi||\Gamma|)\ge|\bar N||S|^2/(\bar\nu^2|\Gamma|)$.

(ii) If $g\in\Sigma$, $p\in\mathcal P$ and $x^{-1}gx\in\widetilde C$, then $x^{-1}gpx=(x^{-1}gx)(x^{-1}px)$ with
$x^{-1}px\in\mathcal P\subset N_0$, and $\widetilde CN_0=\widetilde C$. So $gp$ has at least as many good cosets as $g$.

(iii) Let $\bar N_\ell=\bar N\cap(G_\ell\times1)$, the kernel of the projection of $\bar N$ to the other coordinates.
Its image $\pi_\ell(\bar N_\ell)$ is normal in $\pi_\ell(\bar N)$, so its intersection with the simple group
$\PSL_2(\F_\ell)$ is trivial or everything. If it were trivial, $\PSL_2(\F_\ell)$ would embed in
$\pi_\ell(\bar N)/\pi_\ell(\bar N_\ell)$, which is a quotient of a subgroup of $\prod_{\ell'\ne\ell}G_{\ell'}$; then
$\PSL_2(\F_\ell)$ would be a composition factor of a subgroup of some $\PGL_2(\F_{\ell'})$ with $\ell'\ne\ell$. By
Dickson's classification of the subgroups of $\PSL_2(\F_q)$ \cite[Part~II, Ch.~XII, \S260]{Dic01} (see also
\cite[Kap.~II, Hauptsatz~8.27]{Hup67}), applied to $\PGL_2(\F_{\ell'})\subset\PSL_2(\F_{\ell'^2})$, the nonabelian
composition factors of subgroups of $\PGL_2(\F_{\ell'})$ are
$\PSL_2(\F_{\ell'})$ and $A_5$, and $\PSL_2(\F_\ell)$ is neither when $\ell\ge7$.
\end{proof}

\begin{lemma}[per-factor bound]\label{lem:factor}
For $U\le\PGL_2(\F_\ell)$, put $\nu(U)=\max_y|yU\cap\Inv|$ and $F(U)=|U||\Inv|^2/(\nu(U)^2|G_\ell|)$. If
$U\not\supseteq\PSL_2(\F_\ell)$, then $F(U)\ge F(T_n)=\ell^3/((\ell+1)(\ell^2-1))\ge1-1/\ell$, and $F(U)<1$ only for
$U$ conjugate to $T_n$ (cyclic of order $\ell+1$).
\end{lemma}

\begin{proof}[Proof for $\ell\ge61$]
Put $\ell'=|\Inv|^2/|G_\ell|=\ell^3/(\ell^2-1)>\ell$, so $F(U)=|U|\ell'/\nu(U)^2$. Since $\nu(U)\le|U|$, we have
$F(U)\ge\ell'/|U|$. We first bound $\nu$ for the maximal subgroups. A left coset $yB$ of the Borel subgroup
$B=\mathrm{Stab}(\infty)$ consists of the elements mapping $\infty$ to $P=y\infty$. If $P=\infty$, the involutions in
it fix $\infty$ and are determined by their second fixed point, so there are $\ell$ of them; otherwise they swap
$\infty$ and $P$, and there are $\ell-1$ of them (the elements of $\Nor'\setminus T'$ for the split torus $T'$
fixing $\infty$ and $P$). A left coset of $\Nor(T_s)=\mathrm{Stab}\{0,\infty\}$ consists of the elements mapping
$\{0,\infty\}$ onto a fixed pair $\{P,Q\}$. If $\{P,Q\}=\{0,\infty\}$, it is $\Nor(T_s)$, which contains $\ell$
involutions ($\ell-1$ in $\Nor(T_s)\setminus T_s$ and one in $T_s$). Otherwise one of $P,Q$, say $P$, lies outside
$\{0,\infty\}$, and each involution in the coset swaps $P$ with $0$ or with $\infty$. The involutions swapping $0$ and
$P$ form the set $\Nor'\setminus T'$ for the split torus $T'$ fixing $0$ and $P$, and $T'$ acts simply transitively on
$\PP^1\setminus\{0,P\}$, so at most one of them maps $\infty$ into $\{P,Q\}$; similarly with $0$ and $\infty$
exchanged. So such a coset contains at most $2$ involutions. Hence cosets of $B$ and of $\Nor(T_s)$ contain at most
$\ell$ involutions. A left coset of $\Nor(T_n)$ other than $\Nor(T_n)$ itself contains at most $2$ involutions, since a
M\"obius involution is determined by two swapped pairs; and $\Nor(T_n)$, dihedral of order $2(\ell+1)$, contains
$\ell+2$ involutions (its $\ell+1$ reflections and the involution of $T_n$). Right cosets give the same counts, since
$\iota\mapsto\iota^{-1}$ preserves $\Inv$. We use Dickson's classification \cite[Part~II, Ch.~XII, \S260]{Dic01}, \cite[Kap.~II, Hauptsatz~8.27]{Hup67} (applied
through $\PGL_2(\F_\ell)\subset\PSL_2(\F_{\ell^2})$): every subgroup $U\not\supseteq\PSL_2(\F_\ell)$ lies in $B$,
$\Nor(T_s)$ or $\Nor(T_n)$, or is isomorphic to $A_4$, $S_4$ or $A_5$.
\begin{itemize}
\item $U\subset B$ or $U\subset\Nor(T_s)$: here $\nu\le\min(|U|,\ell)$, so $F\ge|U|\ell'/\min(|U|,\ell)^2\ge\ell'/\ell>1$.
\item $U\subset\Nor(T_n)$ with $|U|<\ell+1$: then $|U|\le2(\ell+1)/3<\ell<\ell'$, so $F\ge\ell'/|U|>1$.
\item $U\subset\Nor(T_n)$ with $|U|=\ell+1$: $U$ has index $2$ in $\Nor(T_n)$, so it is $T_n$ or one of the two
  dihedral subgroups of order $\ell+1$. For $T_n$, $F(T_n)=\ell'/(\ell+1)$ (here $\nu(T_n)=\ell+1$), which is $<1$. A
  dihedral $U$ of order $\ell+1$ contains $(\ell+1)/2$ of the reflections of $\Nor(T_n)$, and the other coset of $U$ in
  $\Nor(T_n)$ contains the other $(\ell+1)/2$; the involution of $T_n$ lies in exactly one of the two. Cosets of $U$ not
  in $\Nor(T_n)$ lie in a coset of $\Nor(T_n)$ other than $\Nor(T_n)$, so contain at most $2$ involutions. Hence
  $\nu(U)\le(\ell+3)/2$ and $F(U)\ge4(\ell+1)\ell'/(\ell+3)^2=4\ell^3/((\ell-1)(\ell+3)^2)>1$.
\item $U=\Nor(T_n)$: $F=2(\ell+1)\ell'/(\ell+2)^2>1$.
\item $U\in\{A_4,S_4,A_5\}$: $F\ge\ell'/60>1$ when $\ell\ge61$.
\end{itemize}
The inequality $\ell^3/((\ell+1)(\ell^2-1))\ge1-1/\ell$ is equivalent to $\ell^4\ge(\ell^2-1)^2$.
\end{proof}

For $5\le\ell\le61$ the lemma was verified by computer: the second implementation (\S\ref{sec:computations}) computed $\nu$
and $F$ for every subgroup of representatives of the maximal classes $B$, $\Nor(T_s)$, $\Nor(T_n)$, $S_4$ and $A_5$
(the last when $\ell\equiv\pm1\bmod5$), for every prime $5\le\ell\le61$. In every case $\min F=F(T_n)$, with $F<1$ only
for $T_n$. The first implementation agrees for $\ell\le13$. For example $\nu=4,5,7,8$ and $F\approx1.95,2.29,2.72,2.86$ for the
dihedral subgroups of order $\ell+1$ at $\ell=5,7,11,13$. So Lemma~\ref{lem:factor}, and hence
Theorem~\ref{thm:B3}, depends on a finite computation for $\ell<61$.

\begin{proof}[Proof of Theorem~\ref{thm:B3}]
Write $\Gamma=\Gamma_P$, $S=S_P$ and $K_0=\ker\pi$, which is central by hypothesis.

\emph{Step 1 (saturation).} Let $Q$ be the set of $\ell\in P$ with $\pi_\ell(\pi(N_0))\supseteq\PSL_2(\F_\ell)$. By
Lemma~\ref{lem:red}(iii), $\pi(N_0)\supseteq P_0:=\prod_{\ell\in Q}\PSL_2(\F_\ell)$. Let $\widehat P=\pi^{-1}(P_0)$, a
central extension of the perfect group $P_0$, and put $\mathcal P=[\widehat P,\widehat P]$. Since $\widehat P=\mathcal
P'K_0'$ for any subgroup $\mathcal P'\le\widehat P$ with $\pi(\mathcal P')=P_0$ (where $K_0'=K_0\cap\widehat P$ is
central), we get $[\widehat P,\widehat P]=[\mathcal P',\mathcal P']$. Taking $\mathcal P'=N_0\cap\widehat P$, which maps
onto $P_0$, shows $\mathcal P\subset N_0$; and $\mathcal P$ is characteristic in $\widehat P\lhd\cG$, so
$\mathcal P\lhd\cG$. Moreover $\pi(\mathcal P)=[P_0,P_0]=P_0$. By Lemma~\ref{lem:red}(ii) the certificate is valid on
$\Sigma:=\widetilde S\mathcal P=\pi^{-1}(SP_0)$, and $SP_0=\prod_{\ell\in Q}G_\ell\times\prod_{\ell\in R}\Inv_\ell$
with $R=P\setminus Q$, because $\Inv_\ell$ meets both cosets of $\PSL_2(\F_\ell)$.

\emph{Step 2 (deleting $Q$).} Let $\varpi:\cG\to\Gamma_R=\prod_{\ell\in R}G_\ell$ be $\pi$ followed by the
projection. Then $\Sigma=\varpi^{-1}(S_R)$, and Lemma~\ref{lem:red}(i) gives, with $\bar N=\varpi(N_0)$,
$\cE\ge|S_R|/|\bar N|$ and $\cA\cE\ge|\bar N||S_R|^2/(\bar\nu^2|\Gamma_R|)$. For $\ell\in R$, $\pi_\ell(\bar N)$ does
not contain $\PSL_2(\F_\ell)$.

\emph{Step 3 (factorisation).} Order $R$ as $\ell_1<\dots<\ell_r$ and let
$U_j=\pi_{\ell_j}(\bar N\cap\prod_{i\ge j}G_{\ell_i})$; each $U_j$ is a subgroup of $\pi_{\ell_j}(\bar N)$, so
$U_j\not\supseteq\PSL_2(\F_{\ell_j})$. Then $|\bar N|=\prod_j|U_j|$, and $\bar\nu\le\prod_j\nu(U_j)$. To see the latter,
project $y\bar N\cap S_R$ to the first coordinate: the image has at most $\nu(U_1)$ points, each fibre is a coset of
$N_1:=\bar N\cap\prod_{i\ge2}G_{\ell_i}$ (placed in the remaining coordinates), and since $S_R$ is a product set it
meets that coset in at most $\nu(N_1)$ points; induction on $r$ finishes. Since
$|S_R|^2/|\Gamma_R|=\prod_j|\Inv_{\ell_j}|^2/|G_{\ell_j}|$, we get $\cA\cE\ge\prod_jF(U_j)$ and
$\cE\ge\prod_j|\Inv_{\ell_j}|/|U_j|$.

\emph{Step 4 (conclusion).} Every subgroup $U\not\supseteq\PSL_2(\F_\ell)$ has $|U|<\ell^2=|\Inv_\ell|$ for $\ell\ge7$
(the maximal ones have orders $\ell(\ell-1)$, $2(\ell\pm1)$, $24$ and, only when $\ell\equiv\pm1\bmod5$, $60$), and $|\Inv_\ell|/|T_n|\ge\ell/2$. Let $J$
be the set of $j$ with $U_j$ conjugate to $T_n$. Then $\cE\ge\prod_{j\in J}\ell_j/2$, and by
Lemma~\ref{lem:factor}, $\cA\cE\ge\prod_{j\in J}(1-1/\ell_j)$. If $\cE>x^{1-\theta}$, the error term alone exceeds
$c\,x\log x$. Otherwise $\prod_{j\in J}\ell_j/2\le x$, so $J$ has $\ll\log x/\log\log x$ elements, all distinct primes,
and by Mertens $\prod_{j\in J}(1-1/\ell_j)\ge\prod_{\ell\le c''\log x}(1-1/\ell)\gg1/\log\log x$. In either case, by the
AM--GM inequality, $B_{\mathrm{cert}}\ge\min\bigl(c\,x\log x,\ 2\sqrt{c\,\cA\cE\Li(x)x^\theta\log x}\bigr)$, which gives
the claim.
\end{proof}

\begin{remark}[on the hypothesis on $\Omega$; arbitrary $\Omega$, a sketch]\label{rem:B3-omega}
The centrality of $\ker\pi$ is used only in Step~1, to find a normal subgroup of $\cG$ inside $N_0$ that maps onto
$P_0$; this is what makes the bound of Theorem~\ref{thm:B3} independent of $P$. If $\Omega$ contains further fields,
for example fields cut out by representations attached to other forms, the kernel may have nonabelian composition
factors. (An earlier version of this note stated Theorem~\ref{thm:B3} for arbitrary $\Omega$; the argument given there
did not justify that.) A short argument, which we give only as a sketch, shows that for \emph{arbitrary} $\Omega$ the
exponent bound survives with a loss depending on $|P|$:
\[
  B_{\mathrm{cert}}(x)\ \ge\ c_1'\,x^{(1+\theta)/2}\,\bigl(\log(|P|+2)\bigr)^{-1/2}\qquad(x\ge x_0),
\]
with $c_1'>0$ depending only on $c$. \emph{Sketch.} Omit Step~1. Since $P$ is independent, $\pi$ is surjective, and
validity on $\widetilde S=\pi^{-1}(S_P)$ is the hypothesis; so Lemma~\ref{lem:red}(i) applies to $\varpi=\pi$ and
$S=S_P$. Factor $\bar N=\pi(N_0)$ over all of $P$ as in Step~3; the relations $|\bar N|=\prod_j|U_j|$ and
$\bar\nu\le\prod_j\nu(U_j)$ do not use $U_j\not\supseteq\PSL_2(\F_{\ell_j})$, so $\cA\cE\ge\prod_jF(U_j)$. Some
$U_j$ may now contain $\PSL_2(\F_{\ell_j})$, but $F(\PGL_2(\F_\ell))=1$ (here $\nu=\ell^2$) and
$F(\PSL_2(\F_\ell))=2\ell^2/(\ell+1)^2>1$ (here $\nu=\ell(\ell+1)/2$, because $D_s$ and $D_n$ lie in different cosets
of $\PSL_2(\F_\ell)$). Together with Lemma~\ref{lem:factor}, $F(U)\ge1$ for every $U\le G_\ell$ not conjugate to
$T_n$, and $F(T_n)\ge1-1/\ell$. Hence $\cA\cE\ge\prod_{\ell\in P}(1-1/\ell)\gg1/\log(|P|+2)$ by Mertens, and the
AM--GM inequality, $B_{\mathrm{cert}}\ge2(c\,\cA\cE\Li(x)x^\theta\log x)^{1/2}$, gives the display. Goursat's lemma
and $\ell\ge7$ are not used. So a certificate over any $\Omega$ can certify the exponent $(1+\theta)/2-\delta$ only if
$\log(|P|+2)\gg_c x^{2\delta}$. Corollary~\ref{cor:B2} covers arbitrary $\Omega$ for single-$\ell$ majorants without
loss.
\end{remark}

\begin{remark}[a better constant]
The nonsplit Cartan certificate ($H=\Nor(T_n)$, $N_0=T_n$, $C$ the nontrivial coset) has $m=(\ell-1)/2$ (split
involutions lie in $(\ell-1)/2$ conjugates of $\Nor(T_n)\setminus T_n$, nonsplit ones in $(\ell+1)/2$), $\cE=\ell$ and
$\cA\cE=\ell/(\ell-1)$. So it improves the constant of \S\ref{sec:A} by a factor $1+2/\ell+O(\ell^{-2})$. It does not
change the exponent.
\end{remark}

\subsection{Per-prime Selberg and large sieves}\label{subsec:B-sieve}

Theorems~\ref{thm:B1} and~\ref{thm:B3} concern a single majorant or a single extension. A sieve combines many primes
$\ell$ through quadratic weights. The finite-group input for $B_{\mathrm{sieve}}$ is a sharp companion of
Theorem~\ref{thm:B1}.

\begin{lemma}\label{lem:S1}
Let $\ell\ge5$ and let $v$ be a real class function on $\PGL_2(\F_\ell)$. If $\mathbb E v=0$, then
$|\mathbb E_{\Inv}v|\le\lambda_A(v)/\ell$. In general,
\[
  \lambda_A(v^2)\ \ge\ \ell\,\mathbb E_{\Inv}v^2-(\ell-1)\,\mathbb E v^2 ,
\]
and this bound is attained, as a bound in terms of the two averages.
\end{lemma}

\begin{proof}
By the proof of Theorem~\ref{thm:B1}, $|\widehat{1_{\Inv}}(\pi)|\le\dim\pi/(\ell^2-1)$ for every $\pi\ne1$. If
$\mathbb Ev=0$, then $\mu_0\,\mathbb E_{\Inv}v=\mathbb E[v1_{\Inv}]=\sum_{\pi\ne1}\widehat v(\pi)\widehat{1_{\Inv}}(\pi)$,
whose absolute value is at most $\lambda_A(v)/(\ell^2-1)$; divide by $\mu_0=\ell/(\ell^2-1)$. For the second bound, the
dual function $\Psi'$ of that proof gives $\lambda_A(h)\ge\mathbb E[h\Psi']=\ell\,\mathbb E_{\Inv}h-(\ell-1)\mathbb E h$
for every real class function $h$; take $h=v^2$. Sharpness is a numerical observation: linear programs, solved in floating
point with HiGHS, for $\ell=5,\dots,97$ give $\min\{\lambda_A(h):h\ge0,\ \mathbb Eh=1,\ \mathbb E_{\Inv}h\ge r\}=1+\ell(r-1)$ for
$r\in[1,\ell-1/\ell]$, with dual solution $\Psi'$.
\end{proof}

Bella\"iche's large sieve \cite[Thm.~14]{Bel16} is the case $v_\ell=(1_{\Inv}-\mu_0)/\sqrt{\mu_0(1-\mu_0)}$ of
Definition~\ref{def:classes}(c): its bound $\Delta/H$ is at least $B_{\mathrm{sieve}}(g/H)$ for
$g(m)=\prod_{\ell\mid m}\sqrt{(1-\mu_0)/\mu_0}$, because $P\equiv1$ on $S_{L^*}$ for this choice and
$g^{\mathsf T}Eg\le(\max_m\sum_nE(m,n))\,H$ for the nonnegative symmetric matrix $E$. Here
$\tilde r_\ell=(\ell^2-\ell-1)/\ell\ge3.8$.

\begin{proposition}[lower bound for $B_{\mathrm{sieve}}$]\label{prop:sieve}
For data as in Definition~\ref{def:classes}(c), suppose $c>0$ and $1\le w(x)\le x^{(1-\theta)/2}$. Then, uniformly in
$L^*$, $(v_\ell)$ and $y$,
\[
  B_{\mathrm{sieve}}(y)\ \ge\ x^{(1+\theta)/2}\,(c\,w(x))^{1/2}\,x^{-O_{\gamma,\eta}(1/\log\log x)}.
\]
In particular no per-prime Selberg sieve in this class gives an exponent below $(1+\theta)/2$. This includes
Bella\"iche's large sieve with any support and any set $L^*$, and its version with induced (Hecke) test functions,
which needs only $\Zt(\theta)$.
\end{proposition}

\begin{proof}
It suffices to take $\lambda=\lambda_A$, since $\lambda_{\mathrm{mon}}(h-\mu(h))\ge\lambda_A(h-\mu(h))$ for every $h$.

\emph{Step 1 (per-prime constants).} Put $\tilde g_\ell=\mathbb E_{\Inv}v_\ell$,
$\sigma_\ell=\lambda_A(v_\ell)/|\tilde g_\ell|$, $\iota_\ell=\lambda_A(v_\ell^2)/\tilde r_\ell$ and
$\rho_\ell=\min(\sigma_\ell,\iota_\ell)$. By Lemma~\ref{lem:S1}, $\sigma_\ell\ge\ell$ and, using
$\mathbb E_{\Inv}v_\ell^2\ge\tilde r_\ell$, $\iota_\ell\ge\ell-(\ell-1)/\tilde r_\ell$. Since $1+\eta\le\tilde
r_\ell\le\mathbb E_{\Inv}v_\ell^2\le1/\mu_0<\ell$, this gives $\rho_\ell\ge1$, $\tilde r_\ell/\rho_\ell\le1.16$, and
$\tilde r_\ell/\rho_\ell^2\le(1+\eta)^2/(\eta^2\ell)$ and $1/\rho_\ell\le(1+\eta)/(\eta\ell)$. For squarefree $k$ write
$\rho(k)=\prod_{\ell\mid k}\rho_\ell$ and $R(k)=\prod_{\ell\mid k}\tilde r_\ell$.

\emph{Step 2 (relaxation).} Averaging $P\ge1$ over the uniform measure on the product set $S_{L^*}$ gives
$\sum_my_m\prod_{\ell\mid m}\tilde g_\ell\ge1$, so $\sum_mz_m\ge1$ for $z_m:=|y_m|\prod_{\ell\mid m}|\tilde g_\ell|$.
Since $\lambda_A$ is multiplicative on $\prod G_\ell$, for $m\ne n$ we have
$|y_my_n|E(m,n)=z_mz_n\prod_{\ell\mid m\triangle n}\sigma_\ell\prod_{\ell\mid\gcd(m,n)}\iota_\ell\ge z_mz_n\rho(\operatorname{lcm}(m,n))$,
and $y_m^2E(m,m)\ge z_m^2(\rho(m)-1/R(m))$. With $Y_1:=\sum_my_m^2=\sum_mz_m^2/R(m)$ and
$Y_2:=\sum_{m,n}z_mz_n\rho(\operatorname{lcm}(m,n))$ this gives
\[
  B_{\mathrm{sieve}}(y)\ \ge\ (\Li(x)-c\,x^\theta w)\,Y_1+c\,x^\theta w\,Y_2.
\]

\emph{Step 3 (a lemma).} If $z\ge0$ is supported on $\omega(m)\le t$, then $(\sum z)^4\le Y_1Y_2\,\Pi\,(2.16)^t$, where
$\Pi=\prod_{\ell\le x^\gamma}(1+\tilde r_\ell/\rho_\ell^2)$; by Step~1, $\Pi\ll(\gamma\log x)^{O_\eta(1)}$. Indeed,
$\rho(\operatorname{lcm}(m,n))\,\rho(\gcd(m,n))=\rho(m)\rho(n)$, so the Cauchy--Schwarz inequality over pairs gives
\[
  \Bigl(\sum_mz_m\Bigr)^4\le Y_2\sum_{m,n}\frac{z_mz_n\,\rho(\gcd(m,n))}{\rho(m)\rho(n)}.
\]
Write $\rho(k)=\sum_{d\mid k}\phi(d)$ with $\phi(d)=\prod_{\ell\mid d}(\rho_\ell-1)\ge0$. The last sum is then
\[
  \sum_d\phi(d)\Bigl(\sum_{d\mid m}\frac{z_m}{\rho(m)}\Bigr)^2.
\]
Cauchy--Schwarz bounds the inner square by $\bigl(\sum_{d\mid m}z_m^2/R(m)\bigr)\,\Pi\,R(d)/\rho(d)^2$, and summing
over $d$ gives at most
\[
  \Pi\sum_m\frac{z_m^2}{R(m)}\prod_{\ell\mid m}\Bigl(1+\frac{\tilde r_\ell}{\rho_\ell}\Bigr)\le\Pi\,(2.16)^tY_1.
\]

\emph{Step 4 (conclusion).} Let $t=\lceil3\log x/\log\log x\rceil$. By Step~1,
$\sum_\ell1/\rho_\ell\ll_\eta\log\log x^\gamma$, so $\sum_{\omega(m)>t}1/\rho(m)\le x^{-3+o(1)}$. Since
$Y_2\ge\sum_mz_m^2\rho(m)$, Cauchy--Schwarz gives $\sum_{\omega(m)>t}z_m\le Y_2^{1/2}x^{-3/2+o(1)}$. If this exceeds
$\frac12$, then $Y_2\ge x^{3-o(1)}$ and $B_{\mathrm{sieve}}(y)\ge c\,x^\theta wY_2$ is far larger than claimed. Otherwise
restrict $z$ to $\omega(m)\le t$; this keeps $\sum z\ge\frac12$ and does not increase $Y_1$ or $Y_2$. Step~3 gives
$Y_1Y_2\ge(16\,\Pi\,2.16^t)^{-1}=x^{-O_{\gamma,\eta}(1/\log\log x)}$. Since $c\,x^\theta w\le\frac12\Li(x)$ for large
$x$, the AM--GM inequality gives $B_{\mathrm{sieve}}(y)\ge2\bigl(\tfrac12\Li(x)\,c\,x^\theta w\,Y_1Y_2\bigr)^{1/2}$, which
is the claim.
\end{proof}

\begin{remark}\label{rem:sieve}
\begin{enumerate}[label=(\roman*),leftmargin=*]
\item \emph{Constants.} For Bella\"iche's two-point functions, $\sigma_\ell>\tilde r_\ell$ and $\iota_\ell>\tilde
  r_\ell$ exactly (a symbolic computation). The following are \emph{numerical observations, not proved}: on the
  divisors of a fixed $m$ the optimal weights appear to be tensor products, which reproduce the majorant
  $1_{\prod_{\ell\mid m}\Inv_\ell}$ of a single compositum and achieve $Y_1Y_2=\prod_{\ell\mid m}\kappa_\ell$ (checked
  for $\ell=5,7,11,13,97$); and in our numerical optimisations, spreading the weights over many primes gained at most a
  factor of order $(\log\log x)^{1/2}$ over a single prime, as in Theorem~\ref{thm:B3}, and never anything in the
  exponent. The lower bound of Proposition~\ref{prop:sieve} is much weaker than these observations, but points the
  same way.
\item \emph{Calibration.} At $\theta=\frac12$ the bound is $x^{3/4-o(1)}$, matching \cite[Thm.~14 and Cor.~15]{Bel16}
  and Serre \cite{Serre81}. So the exponent $2/3$ of \cite{Bel13} cannot be recovered by any choice of support or
  weights in this class.
\item \emph{Not covered.} The proposition does not cover: primes with $\tilde r_\ell<1+\eta$, which barely sieve
  (numerically they never helped); sign-changing constraints $|P|\ge1$; error bounds that exploit cancellation between
  different pairs $(m,n)$, such as $\lambda(F-\mu(F))$ for $F=P^2$ as a whole; non-tensor or multi-function test
  families, such as Zywina's conjugacy sieve; information modulo $\ell^m$; nonlinear selections; and smoothing at the
  level of logarithms. All but the last three produce majorants of $\prod\Inv_\ell$, and so would be covered by
  Conjecture~\ref{conj:product}.
\end{enumerate}
\end{remark}

\begin{conjecture}\label{conj:product}
For distinct primes $\ell_1,\dots,\ell_r\ge5$, every class-function majorant $\varphi$ of
$\prod_i\Inv_{\ell_i}$ on $\prod_iG_{\ell_i}$ satisfies $\mu(\varphi)\lambda_A(\varphi)\ge\prod_i\kappa_{\ell_i}$.
\end{conjecture}
For $\mu(\varphi)\ge1$ this holds because $\lambda_A(\varphi)\ge\mu(\varphi)$. For tensor products $\varphi=\bigotimes_i
\varphi_i$ of nonnegative functions it holds by Theorem~\ref{thm:B1}, which needs no bound on $\mu$, because $\mu$ and
$\lambda_A$ are multiplicative and the $\varphi_i$ can be rescaled to be majorants of the $\Inv_{\ell_i}$. Linear
programs (floating point, HiGHS) over majorants with fixed mean, on grids covering the full range $\mu_0\le\mu\le3$, confirm the conjecture
at the grid points for $(\ell_i)=(5,7)$, $(5,11)$, $(7,11)$, $(11,13)$ and $(5,7,11)$ (first implementation) and for
$(5,7)$ and $(7,13)$ (second implementation); at the grid points with $\mu\ge\frac12$ the minimum of
$\mu\lambda_A$ is $1$. For $\frac12<\mu<1$ no proved statement covers the product case: there the trivial bound
gives only $\mu\lambda_A\ge\mu^2$.

\subsection{Scope and calibration}\label{subsec:B-scope}

\begin{itemize}[leftmargin=*]
\item \emph{Exponent only.} The lower bounds are statements about the functionals of Definition~\ref{def:classes}.
  They are \emph{not} lower bounds for linear explicit-formula methods in general. Zywina's smoothed Chebotarev
  \cite[Thms.~2.2--2.3]{Zyw15} gives $x^{3/4}(\log x)^{-1/2}$ under GRH, which beats the sharp-cutoff form of
  Corollary~\ref{cor:B2} at $\theta=\frac12$ by $(\log x)^{1/2}$. In Definition~\ref{def:classes}(a) this corresponds
  to $w(x)\asymp\log\ell/\log x$ (our interpretation of his method). This $w$ is $<1$ but satisfies
  $w\ge x^{-o(1)}$, as Definition~\ref{def:classes} allows, so Corollary~\ref{cor:B2} applies to it.
\item \emph{Not covered.} The following variants plausibly count as ``linear Chebotarev'' methods. None is known to
  beat the exponent $(1+\theta)/2$, but none is covered by the results above:
  (a) Hecke-shape errors $|C|^{1/2}n_K$ with $|C|>1$, using several $\ell$ at once;
  (b) information modulo $\ell^m$ (the refinement $\tr\equiv0\bmod\ell^m$ is not a full preimage);
  (c) signed or quadratic weights beyond Proposition~\ref{prop:sieve}: sieves with sign-changing constraints, with
  error terms that cancel between pairs, or with non-tensor test functions (Remark~\ref{rem:sieve}(iii));
  (d) nonlinear selections, such as the maximum over $\ell\in[y,2y]$ in Murty--Murty--Saradha's Lemma~4.4, which
  Zywina's proof uses;
  (e) certificates over fields $\Omega$ whose kernel to $\Gamma_P$ is not central, for which we have only the
  sketch in Remark~\ref{rem:B3-omega} (exponent $(1+\theta)/2$ with a loss $(\log(|P|+2))^{-1/2}$).
  Conjecture~\ref{conj:product} would bring (a) and (c) into scope.
\item \emph{Bella\"iche's large sieve.} The arXiv preprint \cite{Bel13} (v1, 2013; the only arXiv version) states,
  under GRH and Artin's conjecture, $x^{2/3+\varepsilon}$ for $a=0$ and $x^{5/7+\varepsilon}$ for $a\ne0$ (non-CM
  elliptic curves), via a large sieve over all $\ell\le Q$. Its proof bounds
  $\prod_{\ell\in m\cup n}\ell=\operatorname{lcm}(m,n)$ by $Q$ for squarefree $m,n\le Q$; for two distinct primes
  $m,n\in(Q/2,Q]$ this product exceeds $Q^2/4$. The published version \cite{Bel16} no longer contains these bounds.
  Its Remark~18(i) and its introduction acknowledge errors in the sieve argument of an earlier version, found by
  Zywina and later by L.~Devin and the referee.\footnote{The v1 exponents are still quoted in the literature under
  the citation of the published paper, e.g.\ \cite[(6)]{CW}, which refers to ``[Be16, Corollaire 17, p.~51]'', the
  numbering and page of arXiv v1.} The corrected \cite[Thm.~14]{Bel16} has exponent
  $(\alpha+4\beta+1)/(2\alpha+4\beta+2)$, where $|G_\ell|/|D_\ell|\sim\ell^\alpha$ and
  $\lambda_A(D_\ell)\ll\ell^\beta$. For $D_\ell=\Inv_\ell$ we have $\alpha=\beta=1$, since
  $|\Inv_\ell|/|G_\ell|=\ell/(\ell^2-1)$ and $\lambda_A(1_{\Inv})=(\ell^2+1)/(\ell+1)$. This gives $3/4$. With
  $x^{1/2}$ replaced by $x^\theta$ the same proof gives $1-(1+\alpha)(1-\theta)/(1+\alpha+2\beta)=(1+\theta)/2$, in
  line with \cite[Rem.~18(iv)]{Bel16}: at $\alpha=1$ the sieve and the single-$\ell$ method tie. The reason is that
  each sieving prime divides the main term by $\asymp\ell$ and multiplies each off-diagonal error by $\asymp\ell$. The
  sieve can also be run without Artin's conjecture, using induced (Hecke) test functions, whose norm
  $\lambda_{\mathrm{mon}}$ is submultiplicative; this needs only $\Zt(\theta)$ for the composita, and again gives
  $(1+\theta)/2$. By Proposition~\ref{prop:sieve}, no choice of support or weights in Bella\"iche's sieve, with or
  without Artin's conjecture, does better.
\item \emph{Bella\"iche's Remark~16.} The closest structural analogue of Theorem~\ref{thm:B1} that we know is
  \cite[Rem.~16]{Bel16}. It concerns the \emph{least-prime} problem: bounding the least $p$ with $\Frob_p\in D$ by means
  of a test function $f$ satisfying his conditions (a) ($f(g)>0$ only for $g\in D$) and (b) ($\sum_gf(g)>0$). He
  shows that if the error term is measured by the norm $\|f\|_2$ of Murty--Murty--Saradha, the best choice is
  always $f=1_D$, whereas with the Littlewood norm other test functions can do better. Theorem~\ref{thm:B1} is the
  corresponding statement for majorants of $\Inv_\ell$ with the Littlewood norm itself. He also shows
  \cite[Thm.~11]{Bel16} that the trace-zero set has Littlewood complexity $\asymp\ell$.
\item \emph{Calibration at $\theta=\frac12$.} The lower bound is $x^{3/4-o(1)}$, the best GRH (and GRH + Artin)
  exponent for $a=0$ in print that we know of: Serre \cite{Serre81}, Zywina \cite{Zyw15} (with $(\log x)^{-1/2}$), and
  \cite[Cor.~15]{Bel16}. The only better exponent we know of is $x^{1/2}\log x$ for non-CM elliptic curves, under GRH,
  Artin's conjecture and a pair correlation conjecture, due to Murty, Murty and Wong as quoted in \cite[(7)]{CW}; pair
  correlation is an input outside the three classes. Elkies' $x^{3/4}$ in weight $2$ is not a Chebotarev argument.
\item \emph{What would be needed.} Granting Artin's conjecture, or Hu's claimed half-plane for Artin $L$-functions,
  does not change the lower bounds: Theorem~\ref{thm:B1} already allows $\lambda_A$. A better exponent would seem to
  need input of a different kind, such as cancellation among the zeros (pair-correlation type), a zero-density
  estimate near $\operatorname{Re}s=\theta$ that is uniform over the Artin $L$-functions of $\PGL_2(\F_\ell)$ as $\ell$
  grows, or non-Galois input such as Elkies' CM lifting.
\end{itemize}

%===============================================================================
\section{Further consequences of \texorpdfstring{$\Zt(\theta)$}{Z(theta)}}\label{sec:further}
%===============================================================================

The three statements of this section follow from $\Zt(\theta)$, or from $\Zt_K(\theta)$ for single fields, by short
arguments or by instantiating published $\theta$-conditional theorems. They are independent of Theorems~A and~B and of
each other. None of them is a substantial new result; we record them, with their status, because they follow from the
same input.

\subsection{Class group torsion}

\begin{proposition}[immediate consequence]\label{prop:C}
Let $\tfrac12\le\theta<1$, let $n\ge2$ and $\ell\ge1$ be integers, and let $\varepsilon>0$. Every number field $K$ of
degree $n$ that satisfies $\Zt_K(\theta)$ has
$|\Cl_K[\ell]|\ll_{n,\ell,\varepsilon,\theta}D_K^{\,1/2-1/(2\ell(n-1))+\varepsilon}$, with an effective implied
constant. In particular, given (I1) or (I2), the bound of Ellenberg--Venkatesh \cite[Prop.~1]{EV}, previously known
under GRH, holds for every number field.
\end{proposition}

\emph{Status.} This is \cite[Lemma~3]{EV} combined with the prime ideal theorem under $\Zt_K(\theta)$. Frei, Sofos
and Widmer (arXiv, 13 October 2026) prove the same bound \cite[Cor.~1.4]{FSW}: for a fixed number field $K$ and an
extension $L/K$ of degree $d$ such that $\zeta_{\widetilde L}$, for the normal closure $\widetilde L$ of $L/K$, satisfies
QRH$(\delta)$, i.e.\ $\Zt_{\widetilde L}(1-\delta)$, $\#\Cl_L[\ell]\ll_{\delta,K,d,\ell,\varepsilon}\Delta_L^{1/2-1/(2\ell(d-1))+\varepsilon}$. It is
\cite[Prop.~3.1]{EV} with their least-splitting-prime bound \cite[Cor.~1.2]{FSW}, deduced from \cite[Thm.~1.1]{FSW}, in
place of Lagarias--Odlyzko under GRH. With $K=\Q$ it contains Proposition~\ref{prop:C}, except that it assumes the
half-plane for $\zeta_{\widetilde L}$ rather than for the field itself (by Aramata--Brauer the former implies the
latter; under (I1) or (I2) both hold). The body of \cite{FSW} states every result conditionally on QRH$(\delta)$, since
its authors ``have currently not verified the arguments'' of OpenAI and Cureton--Thorner (\S1), while its abstract
says that, drawing on those works, they ``make unconditional the Ellenberg--Venkatesh bounds on $\ell$-torsion in
class groups of number fields''. For fields
that are not $\delta$-exceptional in the sense of Pierce, Turnage-Butterbaugh and Wood \cite[Property~1.2]{PTW}
(i.e.\ $\zeta_{\widetilde K}/\zeta$ has no zero in $[1-\delta,1]\times[-(\log D_{\widetilde K})^{2/\delta},(\log
D_{\widetilde K})^{2/\delta}]$, where $\widetilde K$ is the Galois closure), the same conclusion follows from their
effective Chebotarev theorem \cite[Thm.~1.3]{PTW} together with \cite[Lemma~3]{EV}, as carried out in \cite[\S7.2]{PTW}
for the fields in their families. Under $\Zt(31/32)$ no field is $\delta$-exceptional for $\delta<1/32$ (zeros on
$\operatorname{Re}s=31/32$ are not excluded).

\begin{proof}
Ellenberg--Venkatesh \cite[Lemma~3]{EV} state, with $K_0=\Q$ and $d=n$: if $\delta<1/(2\ell(n-1))$ and
$\fp_1,\dots,\fp_M$ are prime ideals of $K$ of norm at most $D_K^\delta$ that are unramified and are not extensions of
prime ideals from any proper subfield of $K$, then $|\Cl_K[\ell]|\ll_{n,\varepsilon,\ell}D_K^{1/2+\varepsilon}M^{-1}$;
the lemma is unconditional. An unramified prime of degree one is never an extension of a prime $\mathfrak q$ of a proper
subfield $K_1$ (it would be the only prime above $\mathfrak q$ and unramified, so of residue degree $[K:K_1]\ge2$);
this is also in \cite[Lemma~2]{HB} and \cite[\S10.1]{LOTZ}. By Proposition~\ref{prop:cheb} with $L=K$ and
$G=\{1\}$, at $x=D_K^\delta$ the number of such primes is
$M\ge D_K^\delta/(\delta\log D_K)(1+o(1))-O_{n,\theta,\delta}(D_K^{\theta\delta}\log D_K)-nD_K^{\delta/2}-n\log D_K
\gg_{n,\delta,\theta}D_K^\delta/\log D_K$ for $D_K\ge D_0(n,\delta,\theta)$ (prime ideals of degree $\ge2$ number at
most $n\sqrt x$, ramified ones at most $n\log D_K$). By Hermite's theorem the finitely many fields with $D_K<D_0$
affect only the constant. Take $\delta=1/(2\ell(n-1))-\varepsilon/2$.
\end{proof}

The bound is uniform in $\theta$, because the primes are needed only at a fixed power level $D_K^\delta$.
Unconditionally, the best known bound for \emph{all} $K$ and $\ell$ is the logarithmic saving of Lemke
Oliver--Zaman \cite[Cor.~1]{LOZ}; Proposition~\ref{prop:C} gives a power saving under (I1) or (I2). Proposition~\ref{prop:C} is not new, or is superseded, for $\ell=2$
\cite{BSTTTZ}; for quadratic and cubic fields with $\ell=3$ \cite{EV,CK}; for $n=2,3$ and the families of the record
\cite{Bha} (\S17); for fields with nilpotent or $\ell$-group Galois closure \cite{Wan,KW}; and for pure cubic and pure
odd-degree fields \cite{HB,Wid}. For non-$D_4$ quartic fields with $\ell=3$ it gives the saving $1/18$; the saving
$1/168$ for quartic fields with Galois group $A_4$ or $S_4$ is attributed to Ellenberg--Venkatesh in
\cite[(1.4)]{EPW} and \cite[\S7]{PTW} (the published \cite[Cor.~1]{EV} states, for quartic fields, only some $\delta>0$). Results for almost
all fields in a family are in \cite{PTW,EPW,LOTZ}. Our survey of these cases is not exhaustive.

\subsection{Elliptic curves with prime order}

\begin{proposition}[instantiation]\label{prop:D}
Let $E/\Q$ be an elliptic curve without CM, of conductor $N_E$, and assume $\Zt(\theta)$ for some $\theta<1$.
\begin{enumerate}[label=(\alph*)]
\item $\#\{p\le x:\ p\nmid N_E,\ |E(\F_p)|\text{ prime}\}\ll_{N_E} x/(\log x)^2$ \textup{(Cojocaru \cite[Prop.~7]{Coj05})}.
\item $\sum_{p\nmid N_E,\ |E(\F_p)|\text{ prime}}1/p<\infty$ \textup{(Cojocaru \cite[Cor.~8]{Coj05})}.
\item If $E$ is not $\Q$-isogenous to a curve with a nontrivial rational torsion point, then for every
  $\varepsilon>0$ and $x\ge x_0(E,\varepsilon)$,
  $\#\{p\le x:\ |E(\F_p)|\text{ prime}\}\le\bigl(\frac5{1-\theta}+\varepsilon\bigr)C_E^{\mathrm{twin}}\frac{x}{(\log x)^2}$
  \textup{(David--Wu \cite[Thm.~1.3]{DW})}, with the constant $C_E^{\mathrm{twin}}$ of \cite{DW}. The factor is
  $160+\varepsilon$ at $\theta=31/32$.
\end{enumerate}
\end{proposition}

\emph{Status.} These are the cited theorems, which are proved under a $\theta$-hypothesis for arbitrary $\theta<1$.
That Cojocaru's results need only a $\theta$-quasi-GRH, and that Zywina's own bound \cite[Thm.~1.3(ii)]{Zyw08}
(constant $22$ under GRH) holds under a $\theta$-quasi-GRH with a $\theta$-dependent constant, is already stated by
Zywina \cite[footnote~1]{Zyw08}; his statements cover number fields and curves with rational torsion on an isogenous
curve as well. The only check needed is that $\Zt(\theta)$ implies the hypotheses used. Cojocaru uses a zero-free
half-plane only for the division fields $\Q(E[d])$, through Lagarias--Odlyzko with error $O(d^3x^\theta\log(dN_Ex))$
\cite[Rem.~12]{Coj05}; these fields vary with $x$, and $\Zt(\theta)$ covers them uniformly. David--Wu's
Hypothesis~3.4 asks for the half-plane for Dedekind zeta functions and for Artin $L$-functions satisfying Artin's
conjecture; in their Thm.~3.8 it is used for Hecke $L$-functions of an abelian quotient, which are covered by
Proposition~\ref{prop:artin} (and by \cite[Cor.~1.2]{CT}). Their level $\xi=2(1-\theta)(1-\varepsilon)/5$ gives the
constant $5/(1-\theta)$. This is an upper bound of the conjectured order, not Koblitz's conjecture. Unconditionally, the best known
bounds are $x/((\log x)\log\log\log x)$ \cite{Coj05} and $(24+o(1))Cx/((\log x)\log\log x)$
\cite[Thm.~1.3(i)]{Zyw08}, and \cite[Rem.~1.4]{Zyw08} lists the convergence of the sum in (b) as open;
Proposition~\ref{prop:D} gives (a)--(c) under (I1) or (I2). Further published
theorems stated under a $\theta$-quasi-GRH, such as the Erd\H{o}s--Kac laws of \cite{MMP} and \cite[Thm.~1, Cor.~2]{CM},
are instantiated in the same way.

\subsection{CM fields with relative class number one}

Let $K$ be a CM field of degree $2n$ with maximal totally real subfield $k$. Let $h^-_K=h_K/h_k$, and put
$I=[1-1/(c\log D_K),1)$ with $c=(2+\sqrt3)/4$.

\begin{proposition}[short deduction from Lee--Kwon]\label{prop:E}
\begin{enumerate}[label=(\alph*)]
\item Assume $\Zt_K(1-1/6000)$. If $n\ge127$ then $h_K^->1$. Hence, given (I1), every CM field with $h^-_K=1$, and in
  particular every CM field with class number one, has degree $\le252$.
\item Let $n\ge109$ and suppose $\zeta_K$ has no real zero in $I$. This holds in particular under $\Zt_K(\theta)$
  whenever $1-\theta>1/(218\,c\,B_{50})\approx1/717$, where $B_{50}=3.524$ is Lee--Kwon's constant; so it holds for
  $\theta=31/32$. Then $h^-_K\ge(1.06136)^n/(5.98607\,n)>1$. Hence, given (I2), every CM field with $h^-_K=1$ has
  degree $\le216$.
\end{enumerate}
Both parts use the numerical constants of Lee--Kwon \cite[Tables~1 and~7]{LK}. Part (a) uses the rows $n=127,128,129$
of \cite[Table~7]{LK} as printed; we have not recomputed them, and with only the rows we recomputed the bound in (a)
becomes $\le258$.
\end{proposition}

\emph{Status.} Lee--Kwon prove $\le216$ for \emph{normal} CM fields \cite[Thm.~1(2)]{LK}, and $\le104$ for all CM
fields under GRH \cite[Thm.~1(3)]{LK}. In their proof of \cite[Prop.~8(2)]{LK} normality is used only in case~(ii), a
simple zero of $\zeta_K$ in $I$ that is not a zero of $\zeta_k$, which normality places in an imaginary quadratic
subfield (via Stark \cite{Stark74} and Bessassi \cite{Bes03}). Proposition~\ref{prop:E} says that their argument extends
to all CM fields once that case is excluded by a zero-free region, and that (I1) or (I2) excludes it sufficiently well.
The method is folklore-level. Unconditionally, the degree bounds for CM fields with $h^-=1$ that we know of are for
normal CM fields \cite{Hof79,Bes03,LK}; Proposition~\ref{prop:E} gives a bound for all CM fields under (I1) or
(I2). The record \cite{Bha} (\S34) removes a Siegel-zero hypothesis for CM fields of
fixed degree.

\begin{proof}[Proof (following {\cite[\S\S3--6]{LK}})]
We use from \cite{LK}, none of which assumes normality: the class number formula \cite[(1.4)]{LK},
$h^-_K\ge2D_K^{1/4}\kappa_K/((2\pi)^n\kappa_k)$, where $\kappa$ is the residue at $s=1$; the upper bound for $\kappa_k$
\cite[Prop.~5]{LK}; the lower bound \cite[Prop.~6(2)]{LK}, valid for totally imaginary $K$ of degree $\ge10$ with root
discriminant $\ge2\pi^2$: $\kappa_K\ge1/(ce^{1/(2c)}\log D_K)$ if $\zeta_K(1-1/(c\log D_K))\le0$, and
$\kappa_K\ge(1-\beta)/(2e^{1/(2c)})$ if $\zeta_K(\beta)\le0$ for some $\beta\in I$; and the discriminant bound
\cite[Lemma~7]{LK}. With the parameters of \cite[Table~7]{LK}, for which \cite[\S6.2]{LK} verifies that the remaining
correction factors are $\ge1$, these give $h^-_K\ge2\kappa_K\,C_7^{\,n}$, where $C_7=C_7(m,b,t)$ is Lee--Kwon's constant,
with $m=n$ for $108\le n\le134$ and $m=134$ for $n>134$.

(b) Under the hypothesis, $\zeta_K<0$ on $I$ (the pole at $1$ has positive residue and there are no zeros). Since
$\zeta_K/\zeta_k$ is entire, $\zeta_k$ has no zero in $I$ either. So only Lee--Kwon's case~(iii) occurs, and their chain
(5.9)$\to$(5.4) with Table~7 gives (1.2) for $n\ge109$. For $2n\ge218$, \cite[Lemma~7(1)]{LK} gives $\log D_K\ge
218B_{50}$, so $1/(c\log D_K)\le1/(218cB_{50})<1-\theta$, and $I$ lies in $\operatorname{Re}s>\theta$.

(a) If $\zeta_K$ has a real zero $\beta\in I$ (Lee--Kwon's cases (i) and (ii)), the second branch of Prop.~6(2) applies
and (I1) gives $1-\beta\ge1/6000$, so $h^-_K\ge C_7^{\,n}/(6000e^{1/(2c)})=C_7^{\,n}/10253.9$. With $C_7(127)=1.0763$
from \cite[Table~7]{LK} this is $1.108>1$ at $n=127$ (still $1.102$ if $C_7$ is lowered by the rounding
$5\cdot10^{-5}$); it is larger for $128\le n\le134$ and increases with $n$ beyond $134$. (At $n=126$ it is $0.938$.)
Otherwise (b) applies. Finally, $h_K=1$ implies $h^-_K=1$, because $h_k\mid h_K$ for CM fields.
\end{proof}

%===============================================================================
\section{Computations}\label{sec:computations}
%===============================================================================

The following finite computations support the text. Except where noted, each was carried out twice, with
independently written code: a first implementation, written with the material, and a second one, written separately
during a review pass (\S\ref{sec:provenance}). All scripts, the script \texttt{run\_all.sh} that runs them, and
their output logs are in the directory \texttt{computations/} of the repository
\url{https://github.com/SlackEight/tau-vanishing}; its \texttt{README.md} maps each script to the statement it checks.
They need Python~3 with NumPy, SciPy (whose HiGHS solver does the linear programs) and SymPy, and \texttt{cypari2}
for the values of $\tau(p)$. All linear programs are solved in floating point with HiGHS; they are not exact rational
computations and are not certified, and no proof in this note rests on their values.
\begin{itemize}[leftmargin=*]
\item Lemma~\ref{lem:group}: brute force over $\PGL_2(\F_\ell)$, $5\le\ell\le23$, and explicit coset enumeration for
  $\ell\in\{5,7,11,13,29\}$. Empirical density of $\tau(p)\equiv0\pmod\ell$ for $p\le10^5$; no $p\le10^5$ has $\tau(p)=0$.
\item Theorem~\ref{thm:B1}: character tables and orthogonality (all primes $\ell\le97$); linear programs
  (floating point, HiGHS, mean fixed exactly) over the full range of $\mu$ (all primes $\ell\le97$ in the first implementation; $\ell\in\{5,7,11,13,17,31,61,97\}$ in
  the second).
\item Remark~\ref{rem:B1-num} (sharpness; one implementation, \texttt{sharpness.py}): the Fourier coefficients of
  $\varphi_1$, their signs, $\lambda_A(\varphi_1)=\ell-1$ and the values of $\Psi$ on nonsplit regular elements,
  symbolically for general $\ell$ (SymPy, both classes of $\ell\bmod4$); the same from an explicit character table for
  every prime $5\le\ell\le199$, together with $\lambda_A$ along the segments between the minimisers (floating point);
  and, without any character table, $\lambda_A(\varphi_1)=\ell-1$ and $\lambda_A(1_{\Inv})=\lambda_0$ as the trace norm
  of the convolution operator on $L^2(\PGL_2(\F_\ell))$ divided by $|G|$, for $\ell\in\{5,7,11,13,17\}$.
\item Conjecture~\ref{conj:product}: linear programs on $G_{\ell_1}\times\dots\times G_{\ell_r}$, over the full
  range of $\mu$, for $(5,7)$, $(5,11)$, $(7,11)$, $(11,13)$, $(5,7,11)$ (first implementation) and $(5,7)$, $(7,13)$
  (second).
\item Lemma~\ref{lem:factor}: all subgroups of the maximal-class representatives, for every prime $5\le\ell\le61$
  (second implementation), and all subgroups for $\ell\le13$ (first).
\item Lemma~\ref{lem:S1} and Proposition~\ref{prop:sieve} (second implementation only): the sharpness linear programs of
  Lemma~\ref{lem:S1} for all primes $5\le\ell\le97$; the two-point identities (symbolic); the tensor-weight constants and
  the numerical sieve optimisations of Remark~\ref{rem:sieve}(i).
\item Proposition~\ref{prop:E}: recomputation of $C_7$ from \cite[Table~1]{LK} and of the $h(n)$ column of
  \cite[Table~7]{LK}; the values $C_7^n/(6000e^{1/(2c)})$ for $126\le n\le134$.
\end{itemize}

%===============================================================================
\section{How this note was prepared, and responsibility}\label{sec:provenance}
%===============================================================================

This note was prepared with substantial assistance from AI systems. The tools and their roles were as follows.
\begin{itemize}[leftmargin=*]
\item \emph{Claude} (Anthropic), large language models run as agents in Anthropic's Claude Code environment, under
  the author's direction: literature searches and reading of sources; proposing and drafting statements and proofs;
  writing the text; writing the Lean code of the formalisation and the Python scripts of \S\ref{sec:computations};
  and separate review passes, each run by an agent that had not written the material it reviewed, with independently
  written code where computations were involved.
\item \emph{AxiomProver} (AxiomMath) generated the proofs of the Cureton--Thorner input (I2), according to its
  authors; it was not used for this note.
\item The \emph{Lean~4 kernel}, \emph{Comparator}, \texttt{lean4lean} and \texttt{nanoda} checked the formal proof
  (Appendix~\ref{sec:lean}); \emph{Python} with NumPy, SciPy/HiGHS, SymPy and PARI carried out the computations.
\end{itemize}
The author chose the problem and the research programme, directed the work, set the verification standards, and
decided what is claimed. No AI system is an author. The author takes responsibility for the content, including the
correctness and novelty claims. The author is verifying the proofs; this version is a preprint, has not been refereed,
and should be read as such. The one formally verified statement is the implication of Theorem~A described in
Appendix~\ref{sec:lean}, under the hypothesis stated there. A summary of the review passes, and of the changes they
led to, is in the file \texttt{AUDITS.md} of the repository.

%===============================================================================
\appendix
\section{The Lean formalisation of Theorem A}\label{sec:lean}
%===============================================================================

The repository \url{https://github.com/SlackEight/tau-vanishing} contains a Lean~4 \cite{Lean4} formalisation (Lean
statements and proofs unchanged since commit \texttt{756d8b4}), built with Lean \texttt{v4.35.0-rc3}, Mathlib
\cite{Mathlib} at \texttt{0493595e}, \cite{CT-Lean} at \texttt{09afbb4c} and TauCeti \cite{TauCeti} at \texttt{c197d7bd}.
Build instructions, requirements and the record of the checks are in the repository's \texttt{README.md}.

\emph{Formal statements.} \texttt{card\_vanishing\_le}: let $\frac12\le\theta<1$ and assume $\Zt(\theta)$ in the form
\texttt{ZeroFreeRight K}~$\theta$ for every number field $K$; let $a:\mathbb N\to\Z$ and $N\ne0$ satisfy
\texttt{LargeProjectiveImage a N} (hypothesis H1). Then for every $\varepsilon>0$ there is $C$ with
$\#\{p\le x \text{ prime}:a(p)=0\}\le C\,x^{(1+\theta)/2+\varepsilon}$ for all real $x\ge2$.
\texttt{card\_vanishing\_le\_ct} is the same with $\theta=31/32$, where $\Zt(31/32)$ is supplied by the theorem
\texttt{DedekindZetaZeroFree.main} of \cite{CT-Lean} (a proved Lean dependency, not a hypothesis).
\texttt{card\_tau\_eq\_zero\_le} and \texttt{card\_discriminant\_\allowbreak coeff\_eq\_zero\_le} are the case $a=\tau$, $N=1$,
under \texttt{LargeProjectiveImage ramanujanTau 1}, where \texttt{ramanujanTau} is defined from $q\prod(1-q^n)^{24}$
and proved equal to the $q$-expansion coefficients of Mathlib's \texttt{ModularForm.discriminant}. The count is the
cardinality of a finite set of primes; $C$ depends on $\theta$, $a$, $N$, $\varepsilon$ and the witness of H1.

\emph{Hypothesis H1.} \texttt{ZeroFreeRight K}~$\theta$ says that $\zeta_K$ (Mathlib's \texttt{NumberField.dedekindZeta})
has a holomorphic continuation to $\C\setminus\{1\}$ with no zeros in $\operatorname{Re}s>\theta$. \texttt{LargeProjectiveImage a N}
says that for every prime $\ell\ge\ell_0$ there are a finite Galois $L/\Q$ and $\rho:\Gal(L/\Q)\to\GL_2(\F_\ell)$ such
that $L$ is unramified at every $p\nmid\ell N$, $\tr\rho(\sigma)=a(p)$ in $\F_\ell$ for every arithmetic Frobenius
$\sigma$ at a prime above such $p$, and the projective image is $\PGL_2(\F_\ell)$. It assumes the existence of the
mod-$\ell$ representations (Deligne \cite{Del71}), their ramification and trace properties, and large projective image
(Ribet \cite{Rib85}, Momose \cite{Mom81}; Swinnerton-Dyer \cite{SwD73} for $\Delta$, $\ell_0=692$); none of this is
formalised. If a reference normalises traces with geometric Frobenius, the dual representation inhabits H1. H1 is
satisfiable independently of Deligne (non-CM elliptic curves, by N\'eron--Ogg--Shafarevich and Serre's open image
theorem \cite{Serre72}), and it fails for $a\equiv0$, as it must.

\emph{Trusted base.} The Lean kernel and toolchain, Mathlib, TauCeti, the repository \cite{CT-Lean} (including its
vendored class field theory library \cite{FAGR}), and five definitions that import only Mathlib
(\texttt{ZeroFreeRight}, \texttt{UnramifiedAt}, \texttt{ResidualRep}, \texttt{LargeProjectiveImage},
\texttt{ramanujanTau}). \texttt{\#print axioms} reports only \texttt{propext}, \texttt{Classical.choice} and
\texttt{Quot.sound}. A related formal development on $\tau$, \cite{AxiomTau}, takes $\tau$ as an axiomatised
function; here $\tau$ is defined and linked to Mathlib's $\Delta$. Comparator \cite{Comparator}, with a challenge file that imports only Mathlib and contains the
five definitions and the four statements, accepted the project with the Lean kernel, and in a second run with the
paranoid kernel, \texttt{lean4lean} \cite{Lean4Lean}, \texttt{nanoda} \cite{nanoda} and \texttt{con-leche}; the
replayed closure includes the proof of \texttt{DedekindZetaZeroFree.main}.

\emph{Differences from Theorem~\ref{thm:A}.} The formal theorem is about any integer sequence satisfying H1, in any
weight; the exponent carries $+\varepsilon$ (the Landau-type route used in Lean loses $x^\varepsilon$); it assumes
$\Zt(\theta)$ for all number fields rather than for the $M_\ell$; the constant is existential; and Theorem~B and
\S\ref{sec:further} are not formalised.

%===============================================================================
\section{What is formally checked: proof dependencies}\label{sec:deps}
%===============================================================================

Table~\ref{tab:deps} lists, for each numbered statement, what its proof uses and how each part is checked. The classes
are: (i) formalised in Lean and checked by the kernel; (ii) a published, refereed result; (iii) an unrefereed preprint;
(iv) a conventional deduction in this note, on paper only and not formalised; (v) a computer verification, where
``exact'' means integer or symbolic arithmetic and ``float'' means floating point (HiGHS linear programs or NumPy).
Script paths are relative to the directory \texttt{computations/} of the repository; Lean names are declaration names
in the directory \texttt{TauVanishing/}, with namespace prefixes omitted. The Lean statements and proofs are unchanged
since commit \texttt{756d8b4} (Appendix~\ref{sec:lean}). $\Zt(\theta)$ (or $\Zt_K(\theta)$), and H1 in Lean, are
hypotheses, not dependencies. Of the two inputs, (I1) is class (iii), and (I2) is class (i) but has no human-readable
proof and is not refereed.

{\small
\begin{longtable}{@{}>{\raggedright\arraybackslash}p{0.17\textwidth}>{\raggedright\arraybackslash}p{0.40\textwidth}>{\raggedright\arraybackslash}p{0.37\textwidth}@{}}
\caption{Dependencies of each statement and how it is checked. Classes: (i) Lean, kernel-checked; (ii) published
refereed result; (iii) unrefereed preprint; (iv) deduction in this note, not formalised; (v) computer verification
(``exact'' or ``float'').}\label{tab:deps}\\
\toprule
item & inputs & class \\
\midrule
\endfirsthead
\toprule
item & inputs & class \\
\midrule
\endhead
\bottomrule
\endfoot
Prop.~\ref{prop:artin} & Brauer induction; Aramata--Brauer; Hecke $L$-functions entire; Frobenius reciprocity
  & (ii) classical; (iv) \\
Prop.~\ref{prop:cheb} & \cite[Thm.~7.1, Lemma~5.4, \S9]{LO}; conductor--discriminant formula; Minkowski's bound
  & (ii); (iv); not in Lean \\
Lemma~\ref{lem:group} & elementary & (iv); (i) \texttt{index\_\allowbreak{}torus}, \texttt{relIndex\_\allowbreak{}torus\_\allowbreak{}torusNormalizer},
  \texttt{sq\_\allowbreak{}eq\_\allowbreak{}neg\_\allowbreak{}det\_\allowbreak{}smul\_\allowbreak{}of\_\allowbreak{}trace\_\allowbreak{}eq\_\allowbreak{}zero}, \texttt{trace\_\allowbreak{}eq\_\allowbreak{}zero\_\allowbreak{}of\_\allowbreak{}sq\_\allowbreak{}scalar},
  \texttt{exists\_\allowbreak{}swapped\_\allowbreak{}cosets} (the bound $\ge(\ell-1)/2$; the exact counts are not formalised); (v) exact,
  \texttt{group\_\allowbreak{}lemma.py}, \texttt{audit/\allowbreak{}cosets.py} \\
Lemma~\ref{lem:dict} & decomposition and inertia theory & (ii); (iv); (i) for the instance used:
  \texttt{card\_\allowbreak{}le\_\allowbreak{}card\_\allowbreak{}inert}, \texttt{mem\_\allowbreak{}zpowers\_\allowbreak{}of\_\allowbreak{}smul\_\allowbreak{}eq} \\
Theorem~\ref{thm:A} & $\Zt(\theta)$ (or $\Zt_{M_\ell}(\theta)$); \cite{Del71}; \cite{Rib85,Mom81}; \cite{SwD73};
  \cite[Prop.~6]{Serre81}; Lemmas~\ref{lem:group}, \ref{lem:dict}; Prop.~\ref{prop:cheb}; Bertrand's postulate
  & sharp form $x^{(1+\theta)/2}$: (iv), with inputs (ii); weaker form $x^{(1+\theta)/2+\varepsilon}$: (i)
  \texttt{card\_\allowbreak{}vanishing\_\allowbreak{}le}, under H1 and $\Zt(\theta)$ for all $K$, by a different route
  (\texttt{primeIdealTheta\_\allowbreak{}uniform}, \texttt{inertTheta\_\allowbreak{}le}, \texttt{log\_\allowbreak{}abs\_\allowbreak{}discr\_\allowbreak{}le} in place of
  Prop.~\ref{prop:cheb} and \cite[Prop.~6]{Serre81}) \\
Cor.~\ref{cor:A-I2} & Thm.~\ref{thm:A}; (I2) & sharp form: (iv) with (i) input; $\varepsilon$-form: (i)
  \texttt{card\_\allowbreak{}vanishing\_\allowbreak{}le\_\allowbreak{}ct}, \texttt{card\_\allowbreak{}tau\_\allowbreak{}eq\_\allowbreak{}zero\_\allowbreak{}le},
  \texttt{card\_\allowbreak{}discriminant\_\allowbreak{}coeff\_\allowbreak{}eq\_\allowbreak{}zero\_\allowbreak{}le} (under H1 for $\tau$) \\
Cor.~\ref{cor:A-I1} & Thm.~\ref{thm:A}; (I1) & (iv) with (iii) input; not in Lean \\
Rem.~\ref{rem:smoothing} & \cite[Thms.~2.2--2.3]{Zyw15} & expectation only; not used \\
Thm.~\ref{thm:B1} & character table of $\PGL_2(\F_\ell)$ \cite[\S5.2]{FH}, \cite{PS83}; dual functions $\Psi$, $\Psi'$, $1$
  & (iv); (ii) for the table; (v) float, \texttt{littlewood.py}, \texttt{audit/\allowbreak{}characters\_\allowbreak{}lp.py} \\
Rem.~\ref{rem:B1-num} & sharpness: Thm.~\ref{thm:B1}; the majorant $\varphi_1$; convexity of $\lambda_A$
  & (iv); (v) exact (symbolic, general $\ell$) and float, \texttt{sharpness.py}; LPs (v) float \\
Cor.~\ref{cor:B2} & Thm.~\ref{thm:B1}; AM--GM; averaging over a normal subgroup; large image & (iv); (ii) \\
Lemma~\ref{lem:conc} & counting & (iv) \\
Lemma~\ref{lem:red} & counting; Goursat; Dickson's classification \cite{Dic01,Hup67} & (iv); (ii) \\
Lemma~\ref{lem:factor} & $\ell\ge61$: Dickson's classification, coset counts; $5\le\ell<61$: enumeration
  & (iv), (ii); (v) exact, \texttt{audit/\allowbreak{}subgroups.py} (every prime $\ell\le61$), \texttt{littlewood.py subgroups}
  ($\ell\le13$) \\
Thm.~\ref{thm:B3} & Lemmas~\ref{lem:conc}--\ref{lem:factor}; central extensions; Mertens; \cite{Rib85,Loe17}
  & (iv); (ii); (v) for $\ell<61$ \\
Rem.~\ref{rem:B3-omega} & sketch & (iv), sketch only \\
Lemma~\ref{lem:S1} & proof of Thm.~\ref{thm:B1} & (iv); sharpness (v) float, \texttt{audit/\allowbreak{}sieve/\allowbreak{}perprime.py} \\
Prop.~\ref{prop:sieve} & Lemma~\ref{lem:S1}; multiplicativity of $\lambda_A$ \cite[Prop.~6]{Bel16};
  Cauchy--Schwarz; Mertens-type bound & (iv); (ii) \\
Rem.~\ref{rem:sieve} & two-point constants; sieve optimisations
  & (v) exact, \texttt{audit/\allowbreak{}sieve/\allowbreak{}twopoint.py}; (v) float, \texttt{tensor.py}, \texttt{sieveqp.py} \\
Conj.~\ref{conj:product} & linear programs & (v) float; evidence only \\
Prop.~\ref{prop:C} & \cite[Lemma~3]{EV}; Prop.~\ref{prop:cheb} ($G=\{1\}$); Hermite
  & (iv); (ii); also \cite[Cor.~1.4]{FSW} (iii) \\
Prop.~\ref{prop:D} & \cite[Prop.~7, Cor.~8]{Coj05}; \cite[Thm.~1.3, Thm.~3.8]{DW}; Prop.~\ref{prop:artin};
  \cite[footnote~1]{Zyw08} & (ii); (iv); (iii) for \cite{Zyw08} \\
Prop.~\ref{prop:E} & \cite{LK}: Props.~5, 6(2), 8(2), Lemma~7, Tables~1 and~7
  & (ii); (iv); (v) float, \texttt{lee\_\allowbreak{}kwon.py}, \texttt{audit/\allowbreak{}lee\_\allowbreak{}kwon\_\allowbreak{}case\_\allowbreak{}z.py}; (a) needs (I1) (iii) \\
\end{longtable}}

%===============================================================================
\begin{thebibliography}{99}
%===============================================================================

\bibitem{AxiomTau}
AxiomMath, \emph{ramanujan-tau-misses-primes} (Lean~4 repository), \url{https://github.com/AxiomMath/ramanujan-tau-misses-primes}.

\bibitem{CT-Lean}
D.~Cureton, J.~Thorner (proofs generated by AxiomProver), \emph{DedekindZetaZeroFree}, Lean~4 formalisation of
\cite{CT}, repository AxiomMath/DedekindZetaZeroFree, \url{https://github.com/AxiomMath/DedekindZetaZeroFree}, commit
\texttt{09afbb4ce6570c1615d58509dc197fc9f0956b7d} (13 October 2026); theorems \texttt{DedekindZetaZeroFree.main} and
\texttt{DedekindZetaZeroFree.heckeCorollary}. Intrinsic identifier
\texttt{swh:1:rev:09afbb4ce6570c1615d58509dc197fc9f0956b7d} (not yet archived as of 14 October 2026). The constant
$31/32$ is not in \cite{CT}~v1.

\bibitem{Bel13}
J.~Bella\"iche, \emph{Th\'eor\`eme de Chebotarev et complexit\'e de Littlewood}, arXiv:1308.1022v1 (2013). (Its sieve bounds
are not in the published version; see \cite[Rem.~18(i)]{Bel16}.)

\bibitem{Bel16}
J.~Bella\"iche, \emph{Th\'eor\`eme de Chebotarev et complexit\'e de Littlewood}, Ann. Sci. \'Ec. Norm. Sup\'er. (4)
\textbf{49} (2016), 579--632, doi:10.24033/asens.2291. Remark numbers are those of the journal version.

\bibitem{Bha}
\emph{Unconditional implications of the Quasi-Riemann Hypothesis}, \texttt{hexagon:2610.00178v1},
\url{https://hexagonmath.org/2610.00178v1}, submitted by S.~Bhattacharya, October 2026. The record states that it was
``Written by GPT~6 Astra'' and its provenance statement says that no human scrutiny was performed.

\bibitem{Bhv}
M.~Bhargava, \emph{Every nonsquare integer not equal to $-1$ is a primitive root modulo a positive lower density of
primes}, \texttt{hexagon:2610.00228v1}, \url{https://hexagonmath.org/2610.00228v1} (14 October 2026). Unrefereed; the
deductions are formalised in Lean with the cited theorems of Cureton--Thorner and OpenAI as hypotheses, according to
the record.

\bibitem{BSTTTZ}
M.~Bhargava, A.~Shankar, T.~Taniguchi, F.~Thorne, J.~Tsimerman, Y.~Zhao, \emph{Bounds on $2$-torsion in class groups of
number fields and integral points on elliptic curves}, J. Amer. Math. Soc. \textbf{33} (2020), 1087--1099.

\bibitem{Bes03}
S.~Bessassi, \emph{Bounds for the degrees of CM-fields of class number one}, Acta Arith. \textbf{106} (2003), no.~3,
213--245.

\bibitem{CK}
S.~Chan, P.~Koymans, \emph{A new pointwise bound for $3$-torsion of class groups}, arXiv:2505.00611.

\bibitem{Coj05}
A.~C.~Cojocaru, \emph{Reductions of an elliptic curve with almost prime orders}, Acta Arith. \textbf{119} (2005),
no.~3, 265--289.

\bibitem{CM}
A.~C.~Cojocaru, M.~Meyer, \emph{Non-CM elliptic curves with infinitely many almost prime Frobenius traces},
J. Number Theory \textbf{252} (2023), 74--117; arXiv:2207.08322.

\bibitem{CW}
A.~C.~Cojocaru, T.~Wang, \emph{Bounds for the distribution of the Frobenius traces associated to a generic abelian
variety}, Math. Z. \textbf{313} (2026), Paper No.~7, doi:10.1007/s00209-026-04024-5; arXiv:2207.02913v2.

\bibitem{CT}
D.~Cureton, J.~Thorner, \emph{A quasi-Riemann hypothesis for Dedekind zeta functions}, \texttt{hexagon:2610.00200v1},
\url{https://hexagonmath.org/2610.00200v1} (12 October 2026). Proofs generated and verified in Lean by AxiomProver,
per the authors; the Lean formalisation is \cite{CT-Lean}.

\bibitem{Del71}
P.~Deligne, \emph{Formes modulaires et repr\'esentations $\ell$-adiques}, S\'eminaire Bourbaki, exp.~355, Lecture
Notes in Math. \textbf{179}, Springer (1971), 139--172.

\bibitem{Dic01}
L.~E.~Dickson, \emph{Linear Groups with an Exposition of the Galois Field Theory}, B.~G.~Teubner, Leipzig, 1901;
Part~II, Chapter~XII (subgroups of $LF(2,p^n)=\PSL_2(\F_{p^n})$), summary in \S260.

\bibitem{DW}
C.~David, J.~Wu, \emph{Almost prime values of the order of elliptic curves over finite fields}, Forum Math.
\textbf{24} (2012), no.~1, 99--119; arXiv:0812.2860.

\bibitem{Elk87}
N.~D.~Elkies, \emph{The existence of infinitely many supersingular primes for every elliptic curve over $\Q$}, Invent.
Math. \textbf{89} (1987), 561--567.

\bibitem{Elk91}
N.~D.~Elkies, \emph{Distribution of supersingular primes}, Ast\'erisque \textbf{198--200} (1991), 127--132.

\bibitem{EV}
J.~S.~Ellenberg, A.~Venkatesh, \emph{Reflection principles and bounds for class group torsion}, Int. Math. Res. Not.
IMRN \textbf{2007}, no.~1, Art. ID rnm002.

\bibitem{EPW}
J.~S.~Ellenberg, L.~B.~Pierce, M.~M.~Wood, \emph{On $\ell$-torsion in class groups of number fields}, Algebra Number
Theory \textbf{11} (2017), no.~8, 1739--1778; arXiv:1606.06103.

\bibitem{FAGR}
\emph{FiniteAbelianGlobalReciprocityv2}, a Lean~4 class field theory library, vendored with compatibility changes in
\cite{CT-Lean} under \texttt{vendor/class-field-theory}; archived copy
\url{https://github.com/PalomarArchive/n-yamaguchi-0729--FiniteAbelianGlobalReciprocityv2--ccc6f4bb14f4}.

\bibitem{FH}
W.~Fulton, J.~Harris, \emph{Representation Theory: A First Course}, Graduate Texts in Mathematics \textbf{129},
Springer, New York, 1991; \S5.2 (representations of $\GL_2(\F_q)$ and $\SL_2(\F_q)$).

\bibitem{FSW}
C.~Frei, E.~Sofos, M.~Widmer, \emph{Some applications of the quasi-Riemann hypothesis to class groups and small
generators of number fields}, arXiv:2610.18491v1 (13 October 2026).

\bibitem{HB}
D.~R.~Heath-Brown, \emph{$\ell$-torsion in class groups via Dirichlet $L$-functions}, arXiv:2412.07701v2 (2025).

\bibitem{Hof79}
J.~Hoffstein, \emph{Some analytic bounds for zeta functions and class numbers}, Invent. Math. \textbf{55} (1979),
37--47.

\bibitem{Hup67}
B.~Huppert, \emph{Endliche Gruppen~I}, Grundlehren der mathematischen Wissenschaften \textbf{134}, Springer,
Berlin, 1967; Kapitel~II, Hauptsatz~8.27.

\bibitem{Hu}
B.~Hu, \emph{A seven-eighths zero-free half-plane for unitary Weil $L$-functions}, \texttt{hexagon:2610.00216v1},
\url{https://hexagonmath.org/2610.00216v1} (13 October 2026).

\bibitem{KW}
J.~Kl\"uners, J.~Wang, \emph{$\ell$-torsion bounds for the class group of number fields with an $\ell$-group as Galois
group}, arXiv:2003.12161.

\bibitem{Comparator}
Lean FRO, \emph{Comparator}, the \texttt{lake comparator} command distributed with Lean~4 \texttt{v4.35.0-rc3}, with
the bundled external checkers \texttt{leanchecker-paranoid}, \texttt{lean4lean}, \texttt{nanoda}, \texttt{con-leche}
and \texttt{con-ron}.

\bibitem{Lean4}
L.~de Moura, S.~Ullrich, \emph{The Lean~4 theorem prover and programming language}, in: Automated Deduction --
CADE~28, Lecture Notes in Comput. Sci. \textbf{12699}, Springer, 2021, 625--635.

\bibitem{Lean4Lean}
M.~Carneiro, \emph{Lean4Lean: Verifying a typechecker for Lean, in Lean}, arXiv:2403.14064.

\bibitem{LO}
J.~C.~Lagarias, A.~M.~Odlyzko, \emph{Effective versions of the Chebotarev density theorem}, in: Algebraic Number Fields
(Durham, 1975), Academic Press, London, 1977, 409--464.

\bibitem{LK}
G.-N.~Lee, S.-H.~Kwon, \emph{CM-fields with relative class number one}, Math. Comp. \textbf{75} (2006), no.~254,
997--1013.

\bibitem{Leh47}
D.~H.~Lehmer, \emph{The vanishing of Ramanujan's function $\tau(n)$}, Duke Math. J. \textbf{14} (1947), 429--433.

\bibitem{LOTZ}
R.~J.~Lemke Oliver, J.~Thorner, A.~Zaman, \emph{An approximate form of Artin's holomorphy conjecture and non-vanishing
of Artin $L$-functions}, Invent. Math. \textbf{235} (2024), 893--971; arXiv:2012.14422.

\bibitem{LOZ}
R.~J.~Lemke Oliver, A.~Zaman, \emph{Improving the trivial bound for $\ell$-torsion in class groups}, arXiv:2502.03464.

\bibitem{Loe17}
D.~Loeffler, \emph{Images of adelic Galois representations for modular forms}, Glasg. Math. J. \textbf{59} (2017),
11--25.

\bibitem{Mathlib}
The mathlib Community, \emph{The Lean mathematical library}, in: Proceedings of the 9th ACM SIGPLAN International
Conference on Certified Programs and Proofs (CPP 2020), ACM, 2020, 367--381.

\bibitem{Mom81}
F.~Momose, \emph{On the $\ell$-adic representations attached to modular forms}, J. Fac. Sci. Univ. Tokyo Sect. IA
Math. \textbf{28} (1981), 89--109.

\bibitem{Murphy}
\emph{riemann-hypothesis-lean} (GitHub repository DrMurphyIsIn/riemann-hypothesis-lean), \texttt{README.md}, \S6,
commit \texttt{21762fe7cef795014021507f40eaa4918b6b8580} (13 October 2026),
\url{https://github.com/DrMurphyIsIn/riemann-hypothesis-lean}.

\bibitem{MMP}
M.~R.~Murty, V.~K.~Murty, S.~Pujahari, \emph{An all-purpose Erd\H{o}s--Kac theorem}, Math. Z. \textbf{305} (2023),
Paper No.~45, doi:10.1007/s00209-023-03370-y.

\bibitem{MMS}
M.~R.~Murty, V.~K.~Murty, N.~Saradha, \emph{Modular forms and the Chebotarev density theorem}, Amer. J. Math.
\textbf{110} (1988), 253--281.

\bibitem{nanoda}
C.~Bailey, \emph{nanoda\_lib}, a type checker for Lean~4, \url{https://github.com/ammkrn/nanoda_lib}.

\bibitem{OAI-QRH}
OpenAI, \emph{The Quasi-Riemann Hypothesis: a zero-free half-plane $\operatorname{Re}(s)>7/8$}, preprint,
30 September 2026, \url{https://github.com/openai/math}
(\texttt{preprints/The-Quasi-Riemann-Hypothesis-September-30-2026}).

\bibitem{OAI-Lean}
OpenAI, \texttt{openai/math} Lean formalisation, commit \texttt{fd4aeeb2ee4fc729c18d98444fed42fd0529eeeb}
(7 October 2026); theorem \texttt{OAI.riemannZeta\_ne\_zero\_of\_seven\_eighths\_lt\_re}.

\bibitem{OAI-PR}
OpenAI, \emph{Primitive roots for every admissible integer base}, preprint, 4 October 2026,
\url{https://github.com/openai/math}.

\bibitem{PS83}
I.~Piatetski-Shapiro, \emph{Complex Representations of $\GL(2,K)$ for Finite Fields $K$}, Contemporary Mathematics
\textbf{16}, American Mathematical Society, Providence, RI, 1983.

\bibitem{PTW}
L.~B.~Pierce, C.~L.~Turnage-Butterbaugh, M.~M.~Wood, \emph{An effective Chebotarev density theorem for families of
number fields, with an application to $\ell$-torsion in class groups}, Invent. Math. \textbf{219} (2020), 701--778;
arXiv:1709.09637.

\bibitem{Rib85}
K.~A.~Ribet, \emph{On $\ell$-adic representations attached to modular forms. II}, Glasgow Math. J. \textbf{27} (1985),
185--194.

\bibitem{Serre81}
J.-P.~Serre, \emph{Quelques applications du th\'eor\`eme de densit\'e de Chebotarev}, Publ. Math. IH\'ES \textbf{54}
(1981), 123--201.

\bibitem{Serre72}
J.-P.~Serre, \emph{Propri\'et\'es galoisiennes des points d'ordre fini des courbes elliptiques}, Invent. Math.
\textbf{15} (1972), 259--331.

\bibitem{Stark74}
H.~M.~Stark, \emph{Some effective cases of the Brauer--Siegel theorem}, Invent. Math. \textbf{23} (1974), 135--152.

\bibitem{SwD73}
H.~P.~F.~Swinnerton-Dyer, \emph{On $\ell$-adic representations and congruences for coefficients of modular forms}, in:
Modular Functions of One Variable III, Lecture Notes in Math. \textbf{350}, Springer, 1973, 1--55.

\bibitem{TauCeti}
The Tau Ceti project, \emph{TauCeti} (Lean~4 library), \url{https://github.com/TauCetiProject/TauCeti}, commit
\texttt{c197d7bdbae257e85d23443c89aa7cd34c13410b}.

\bibitem{TZ}
J.~Thorner, A.~Zaman, \emph{A Chebotarev variant of the Brun--Titchmarsh theorem and bounds for the Lang--Trotter
conjectures}, Int. Math. Res. Not. IMRN \textbf{2018}, no.~16, 4991--5027; arXiv:1606.09238.

\bibitem{Wan}
J.~Wang, \emph{Pointwise bound for $\ell$-torsion in class groups II: nilpotent extensions}, arXiv:2006.10295.

\bibitem{Wid}
M.~Widmer, \emph{A short remark on the $\ell$-torsion part of class groups}, Math. Proc. Cambridge Philos. Soc.,
published online 30 July 2026; arXiv:2603.08121.

\bibitem{Zyw08}
D.~Zywina, \emph{The large sieve and Galois representations}, arXiv:0812.2222 (2008). (We found no journal version.)

\bibitem{Zyw15}
D.~Zywina, \emph{Bounds for the Lang--Trotter conjectures}, in: SCHOLAR---a scientific celebration highlighting open
lines of arithmetic research, Contemp. Math. \textbf{655}, Amer. Math. Soc., 2015, 235--256; arXiv:1508.07682.

\end{thebibliography}

\end{document}
